\input{preamble}

% Bien, commençons ici.
%
\begin{document}

\title{Modules adéquats}


\maketitle

\phantomsection
\label{section-phantom}

\tableofcontents

\section{Introduction}
\label{section-introduction}

\noindent
Pour tout schéma $X$, la catégorie $\QCoh(\mathcal{O}_X)$
des modules quasi-cohérents est abélienne et constitue une sous-catégorie
de Serre faible de la catégorie abélienne de tous les $\mathcal{O}_X$-modules.
Il en va de même de la catégorie des modules quasi-cohérents sur
un espace algébrique $X$, considérée comme sous-catégorie de la catégorie
de tous les $\mathcal{O}_X$-modules sur le petit site étale de $X$.
De plus, pour un morphisme quasi-compact et quasi-séparé
$f : X \to Y$, l'image directe $f_*$ et les images directes supérieures
préservent la quasi-cohérence.

\medskip\noindent
Soient maintenant $X$ un schéma et $\mathcal{O}$ le faisceau structural
sur l'un des grands sites de $X$, par exemple le grand site fppf.
La catégorie des $\mathcal{O}$-modules quasi-cohérents est abélienne
(elle est en fait équivalente à la catégorie usuelle des
$\mathcal{O}_X$-modules quasi-cohérents sur le schéma $X$, mentionnée ci-dessus),
mais son plongement dans $\textit{Mod}(\mathcal{O})$ n'est pas exact.
Un exemple est l'application de modules quasi-cohérents
$$
\mathcal{O}_{\mathbf{A}^1_k}
\longrightarrow
\mathcal{O}_{\mathbf{A}^1_k}
$$
sur $\mathbf{A}^1_k = \Spec(k[x])$ donnée par la multiplication par $x$.
Dans la catégorie abélienne des faisceaux quasi-cohérents, cette application
est injective, tandis que dans la catégorie abélienne de tous les
$\mathcal{O}$-modules sur le grand site de $\mathbf{A}^1_k$, elle a un noyau
non nul, comme on le voit en l'évaluant sur les sections au-dessus de
$\Spec(k[x]/(x)) = \Spec(k)$.
De plus, pour un morphisme quasi-compact et quasi-séparé
$f : X \to Y$, le foncteur $f_{big, *}$ ne préserve pas la quasi-cohérence.

\medskip\noindent
Dans ce chapitre, nous introduisons la catégorie de ce que nous appellerons
les modules adéquats, étroitement liée à celle des modules quasi-cohérents,
qui remédie aux deux difficultés mentionnées ci-dessus. Une autre solution,
que nous mettrons en œuvre lorsque nous traiterons des modules quasi-cohérents
sur les champs algébriques, consiste à considérer les $\mathcal{O}$-modules
localement quasi-cohérents et satisfaisant à la propriété de changement de base plat.
Voir Cohomologie des champs, section
\ref{stacks-cohomology-section-loc-qcoh-flat-base-change},
Cohomologie des champs, remarque
\ref{stacks-cohomology-remark-bousfield-colocalization}, et
Catégories dérivées des champs, section
\ref{stacks-perfect-section-derived}.






\section{Conventions}
\label{section-conventions}

\noindent
Dans ce chapitre, nous fixons
$\tau \in \{Zar, \etale, smooth, syntomic, fppf\}$
ainsi qu'un grand $\tau$-site $\Sch_\tau$ comme dans
Topologies, section \ref{topologies-section-procedure}.
Tous les schémas seront des objets de $\Sch_\tau$.
En particulier, étant donné un schéma $S$, nous obtenons les sites
$(\textit{Aff}/S)_\tau \subset (\Sch/S)_\tau$.
Le faisceau structural $\mathcal{O}$ sur ces sites est défini par
la règle $\mathcal{O}(T) = \Gamma(T, \mathcal{O}_T)$.

\medskip\noindent
Tous les anneaux $A$ seront tels que $\Spec(A)$ soit (isomorphe à) un
objet de $\Sch_\tau$. Étant donné un anneau $A$, nous notons
$\textit{Alg}_A$ la catégorie de $A$-algèbres dont les objets sont les
$A$-algèbres $B$ de la forme $B = \Gamma(U, \mathcal{O}_U)$,
où $U$ est un objet affine de $\Sch_\tau$. Ainsi, étant donné un
schéma affine $S = \Spec(A)$, le foncteur
$$
(\textit{Aff}/S)_\tau \longrightarrow \textit{Alg}_A,
\quad
U \longmapsto \mathcal{O}(U)
$$
est une équivalence.





\section{Foncteurs adéquats}
\label{section-quasi-coherent}

\noindent
Dans cette section, nous étudions un sujet étroitement lié aux
images directes des faisceaux quasi-cohérents. La majeure partie de ces
résultats est tirée de l'article \cite{Jaffe}.

\begin{definition}
\label{definition-module-valued-functor}
Soit $A$ un anneau. Un {\it foncteur à valeurs dans les modules} est un foncteur
$F : \textit{Alg}_A \to \textit{Ab}$ tel que
\begin{enumerate}
\item pour tout objet $B$ de $\textit{Alg}_A$, le groupe
$F(B)$ est muni d'une structure de $B$-module, et
\item pour tout morphisme $B \to B'$ de $\textit{Alg}_A$, l'application
$F(B) \to F(B')$ est $B$-linéaire.
\end{enumerate}
Un {\it morphisme de foncteurs à valeurs dans les modules} est une transformation
de foncteurs $\varphi : F \to G$ telle que $F(B) \to G(B)$ soit $B$-linéaire
pour tout $B \in \Ob(\textit{Alg}_A)$.
\end{definition}

\noindent
Soit $S = \Spec(A)$ un schéma affine.
La catégorie des foncteurs à valeurs dans les modules sur $\textit{Alg}_A$ est
équivalente à la catégorie
$\textit{PMod}((\textit{Aff}/S)_\tau, \mathcal{O})$
des préfaisceaux de $\mathcal{O}$-modules. L'équivalence associe
au foncteur à valeurs dans les modules $F$ le
préfaisceau $\mathcal{F}$ défini par la règle
$\mathcal{F}(U) = F(\mathcal{O}(U))$.
Cela résulte de l'équivalence
$(\textit{Aff}/S)_\tau \to \textit{Alg}_A$, $U \mapsto \mathcal{O}(U)$,
donnée dans la section \ref{section-conventions}.
Le quasi-inverse est défini par $F(B) = \mathcal{F}(\Spec(B))$.

\medskip\noindent
Un cas particulier important de foncteur à valeurs dans les modules s'obtient ainsi.
Soit $M$ un $A$-module. Nous noterons $\underline{M}$ le
foncteur à valeurs dans les modules $B \mapsto M \otimes_A B$ (muni de la structure
évidente de $B$-module). Notons que si $M \to N$ est une application de $A$-modules,
il lui est associé un morphisme $\underline{M} \to \underline{N}$ de foncteurs
à valeurs dans les modules. Réciproquement, tout morphisme de tels foncteurs
$\underline{M} \to \underline{N}$ provient d'une application de $A$-modules $M \to N$,
comme le lecteur le verra en évaluant en $B = A$. Autrement dit,
$\text{Mod}_A$ est une sous-catégorie pleine
de la catégorie des foncteurs à valeurs dans les modules sur $\textit{Alg}_A$.

\medskip\noindent
Étant donnée une application de $A$-modules $\varphi : M \to N$, on a
$\Coker(\underline{M} \to \underline{N}) =
\underline{Q}$, où $Q = \Coker(M \to N)$, car $\otimes$
est exact à droite. Mais il n'en va pas de même
du noyau en général : par exemple, une application injective de
$A$-modules peut cesser de l'être après changement de base. La définition
suivante est donc naturelle.

\begin{definition}
\label{definition-adequate-functor}
Soit $A$ un anneau. Un foncteur à valeurs dans les modules $F$ sur $\textit{Alg}_A$ est
dit
\begin{enumerate}
\item {\it adéquat} s'il existe une
application de $A$-modules $M \to N$ telle que $F$ soit isomorphe à
$\Ker(\underline{M} \to \underline{N})$.
\item {\it linéairement adéquat} si $F$ est isomorphe au
noyau d'une application $\underline{A^{\oplus n}} \to \underline{A^{\oplus m}}$.
\end{enumerate}
\end{definition}

\noindent
Notons que $F$ est adéquat si et seulement s'il existe une
suite exacte $0 \to F \to \underline{M} \to \underline{N}$, et que
$F$ est linéairement adéquat si et seulement s'il existe une suite exacte
$0 \to F \to \underline{A^{\oplus n}} \to \underline{A^{\oplus m}}$.

\medskip\noindent
Soit $A$ un anneau. Dans cette section, nous montrerons que la catégorie des foncteurs
adéquats sur $\textit{Alg}_A$ est abélienne
(lemmes \ref{lemma-cokernel-adequate} et \ref{lemma-kernel-adequate})
et possède un ensemble de générateurs
(lemme \ref{lemma-adequate-surjection-from-linear}).
Nous verrons aussi que c'est une sous-catégorie de Serre faible de la catégorie
de tous les foncteurs à valeurs dans les modules sur $\textit{Alg}_A$
(lemme \ref{lemma-extension-adequate})
et qu'elle admet des limites inductives arbitraires
(lemme \ref{lemma-colimit-adequate}).

\begin{lemma}
\label{lemma-adequate-finite-presentation}
Soit $A$ un anneau.
Soit $F$ un foncteur adéquat sur $\textit{Alg}_A$.
Si $B = \colim B_i$ est une limite inductive
filtrante de $A$-algèbres, alors $F(B) = \colim F(B_i)$.
\end{lemma}

\begin{proof}
Cela résulte de ce que, pour tout $A$-module $M$, on a
$M \otimes_A B = \colim M \otimes_A B_i$ (voir
Algèbre, lemme \ref{algebra-lemma-tensor-products-commute-with-limits})
et de ce que les limites inductives filtrantes commutent aux suites exactes ; voir
Algèbre, lemme \ref{algebra-lemma-directed-colimit-exact}.
\end{proof}

\begin{remark}
\label{remark-settheoretic}
Considérons la catégorie $\textit{Alg}_{fp, A}$ dont les objets sont les $A$-algèbres
$B$ de la forme $B = A[x_1, \ldots, x_n]/(f_1, \ldots, f_m)$ et dont les
morphismes sont les homomorphismes de $A$-algèbres. Toute $A$-algèbre $B$ est une limite
inductive filtrante de $A$-algèbres de présentation finie, c'est-à-dire une limite
inductive filtrante d'objets de $\textit{Alg}_{fp, A}$. Le
lemme \ref{lemma-adequate-finite-presentation}
montre donc que tout foncteur adéquat $F$ est déterminé par sa restriction à
$\textit{Alg}_{fp, A}$. Pour certaines questions, on peut par conséquent se restreindre
aux foncteurs sur $\textit{Alg}_{fp, A}$. Par exemple, la catégorie des foncteurs
adéquats ne dépend pas du choix du grand $\tau$-site
effectué dans la
section \ref{section-conventions}.
\end{remark}

\begin{lemma}
\label{lemma-adequate-flat}
Soit $A$ un anneau.
Soit $F$ un foncteur adéquat sur $\textit{Alg}_A$.
Si $B \to B'$ est plat, alors $F(B) \otimes_B B' \to F(B')$
est un isomorphisme.
\end{lemma}

\begin{proof}
Choisissons une suite exacte $0 \to F \to \underline{M} \to \underline{N}$.
Elle donne le diagramme
$$
\xymatrix{
0 \ar[r] & F(B) \otimes_B B' \ar[r] \ar[d] &
(M \otimes_A B)\otimes_B B' \ar[r] \ar[d] &
(N \otimes_A B)\otimes_B B' \ar[d] \\
0 \ar[r] & F(B') \ar[r] &
M \otimes_A B' \ar[r] &
N \otimes_A B'
}
$$
dont les lignes sont exactes (celle du haut parce que $B \to B'$ est plat).
Puisque les deux flèches verticales de droite sont des isomorphismes, celle
de gauche en est également un.
\end{proof}

\begin{lemma}
\label{lemma-adequate-surjection-from-linear}
Soit $A$ un anneau.
Soit $F$ un foncteur adéquat sur $\textit{Alg}_A$. Il existe alors une
surjection $L \to F$, où $L$ est une somme directe de foncteurs linéairement adéquats.
\end{lemma}

\begin{proof}
Choisissons une suite exacte $0 \to F \to \underline{M} \to \underline{N}$,
où $\underline{M} \to \underline{N}$ est donnée par
$\varphi : M \to N$. D'après le
lemme \ref{lemma-adequate-finite-presentation},
il suffit de construire $L \to F$ de sorte que $L(B) \to F(B)$ soit surjective
pour toute $A$-algèbre de présentation finie $B$. Il suffit donc, étant donnés
une $A$-algèbre de présentation finie $B$ et un élément $\xi \in F(B)$,
de construire une application $L \to F$, avec $L$ linéairement adéquat, telle que $\xi$
appartienne à l'image de $L(B) \to F(B)$.
(En effet, les classes d'isomorphisme de tels couples $(B, \xi)$ forment un ensemble.)

\medskip\noindent
Pour cela, notons $\sum_{i = 1, \ldots, n} m_i \otimes b_i$ l'image de
$\xi$ dans $\underline{M}(B) = M \otimes_A B$. Nous savons que
$\sum \varphi(m_i) \otimes b_i = 0$ dans $N \otimes_A B$.
Comme $N$ est une limite inductive filtrante de $A$-modules de présentation finie,
on peut trouver un $A$-module de présentation finie $N'$ et un diagramme commutatif
de $A$-modules
$$
\xymatrix{
A^{\oplus n} \ar[r] \ar[d]_{m_1, \ldots, m_n} & N' \ar[d] \\
M \ar[r] & N
}
$$
tels que $(b_1, \ldots, b_n)$ ait une image nulle dans $N' \otimes_A B$.
Choisissons une présentation $A^{\oplus l} \to A^{\oplus k} \to N' \to 0$.
Choisissons un relèvement $A^{\oplus n} \to A^{\oplus k}$ de l'application
$A^{\oplus n} \to N'$ du diagramme. On voit alors qu'il existe
$(c_1, \ldots, c_l) \in B^{\oplus l}$ tel que
$(b_1, \ldots, b_n, c_1, \ldots, c_l)$ ait une image nulle dans $B^{\oplus k}$
par l'application $B^{\oplus n} \oplus B^{\oplus l} \to B^{\oplus k}$.
Considérons le diagramme commutatif
$$
\xymatrix{
A^{\oplus n} \oplus A^{\oplus l} \ar[r] \ar[d] & A^{\oplus k} \ar[d] \\
M \ar[r] & N
}
$$
où la flèche verticale de gauche est nulle sur le facteur direct $A^{\oplus l}$.
On voit alors que le noyau $L$ de $\underline{A^{\oplus n + l}}
\to \underline{A^{\oplus k}}$ convient, car l'élément
$(b_1, \ldots, b_n, c_1, \ldots, c_l) \in L(B)$ a pour image $\xi$.
\end{proof}

\noindent
Considérons une $A$-algèbre graduée $B = \bigoplus_{d \geq 0} B_d$. Il existe alors
deux homomorphismes de $A$-algèbres $p, a : B \to B[t, t^{-1}]$, à savoir
$p : b \mapsto b$ et $a : b \mapsto t^{\deg(b)} b$, où $b$ est homogène.
Si $F$ est un foncteur à valeurs dans les modules sur $\textit{Alg}_A$, posons
\begin{equation}
\label{equation-weight-k}
F(B)^{(k)} = \{\xi \in F(B) \mid t^k F(p)(\xi) = F(a)(\xi)\}.
\end{equation}
Pour les foncteurs qui se comportent bien relativement aux extensions plates
d'anneaux, cela donne une décomposition en somme directe. Cela revient au fait
que les représentations de $\mathbf{G}_m$ sont complètement réductibles.

\begin{lemma}
\label{lemma-flat-functor-split}
Soit $A$ un anneau.
Soit $F$ un foncteur à valeurs dans les modules sur $\textit{Alg}_A$.
Supposons que, pour $B \to B'$ plat, l'application
$F(B) \otimes_B B' \to F(B')$ soit un isomorphisme.
Soit $B$ une $A$-algèbre graduée. Alors
\begin{enumerate}
\item $F(B) = \bigoplus_{k \in \mathbf{Z}} F(B)^{(k)}$, et
\item l'application $B \to B_0 \to B$ induit une application $F(B) \to F(B)$
dont l'image est contenue dans $F(B)^{(0)}$.
\end{enumerate}
\end{lemma}

\begin{proof}
Soit $x \in F(B)$. L'application $p : B \to B[t, t^{-1}]$ est libre,
donc nous savons que
$$
F(B[t, t^{-1}]) =
\bigoplus\nolimits_{k \in \mathbf{Z}} F(p)(F(B)) \cdot t^k =
\bigoplus\nolimits_{k \in \mathbf{Z}} F(B) \cdot t^k
$$
comme indiqué, nous omettons $F(p)$ dans la suite de la démonstration.
Écrivons $F(a)(x) = \sum t^k x_k$, avec $x_k \in F(B)$.
Notons $\epsilon : B[t, t^{-1}] \to B$
l'homomorphisme de $B$-algèbres $t \mapsto 1$. Notons que les composés
$\epsilon \circ p, \epsilon \circ a : B \to B[t, t^{-1}] \to B$ sont
l'identité. On a donc
$$
x = F(\epsilon)(F(a)(x)) = F(\epsilon)(\sum t^k x_k) = \sum x_k.
$$
D'autre part, nous affirmons que $x_k \in F(B)^{(k)}$. Considérons en effet
le diagramme commutatif
$$
\xymatrix{
B \ar[r]_a \ar[d]_{a'} &
B[t, t^{-1}] \ar[d]^f \\
B[s, s^{-1}] \ar[r]^-g &
B[t, s, t^{-1}, s^{-1}]
}
$$
où $a'(b) = s^{\deg(b)}b$, $f(b) = b$, $f(t) = st$,
$g(b) = t^{\deg(b)}b$ et $g(s) = s$. Alors
$$
F(g)(F(a'))(x) = F(g)(\sum s^k x_k) =
\sum s^k F(a)(x_k)
$$
et, en suivant l'autre chemin, on obtient
$$
F(f)(F(a))(x) = F(f)(\sum t^k x_k) = \sum (st)^k x_k.
$$
{\emergencystretch=2em
Puisque $B \to B[s, t, s^{-1}, t^{-1}]$ est libre, on a
$F(B[t, s, t^{-1}, s^{-1}]) =
\bigoplus_{k, l \in \mathbf{Z}} F(B) \cdot t^ks^l$ ;
en comparant les coefficients dans les expressions ci-dessus, on trouve
$F(a)(x_k) = t^k x_k$, comme voulu.\par}

\medskip\noindent
Enfin, l'image de $F(B_0) \to F(B)$ est contenue dans $F(B)^{(0)}$
car $B_0 \to B \xrightarrow{a} B[t, t^{-1}]$ est égal à
$B_0 \to B \xrightarrow{p} B[t, t^{-1}]$.
\end{proof}

\noindent
Comme cas particulier du
lemme \ref{lemma-flat-functor-split},
notons que
$$
\underline{M}(B)^{(k)} = M \otimes_A B_k
$$
où $B_k$ est la composante de degré $k$ de la $A$-algèbre graduée $B$.

\begin{lemma}
\label{lemma-lift-map}
Soit $A$ un anneau. Étant donné un diagramme en traits pleins
$$
\xymatrix{
0 \ar[r] &
L \ar[d]_\varphi \ar[r] &
\underline{A^{\oplus n}} \ar[r] \ar@{..>}[ld] &
\underline{A^{\oplus m}} \\
& \underline{M}
}
$$
de foncteurs à valeurs dans les modules sur $\textit{Alg}_A$
dont la ligne est exacte, il existe une flèche pointillée rendant le diagramme commutatif.
\end{lemma}

\begin{proof}
Supposons que l'application $A^{\oplus n} \to A^{\oplus m}$ soit donnée par la
matrice de taille $m \times n$ notée $(a_{ij})$. Considérons l'anneau
$B = A[x_1, \ldots, x_n]/(\sum a_{ij}x_j)$. L'élément
$(x_1, \ldots, x_n) \in \underline{A^{\oplus n}}(B)$ a une image nulle dans
$\underline{A^{\oplus m}}(B)$ ; il est donc l'image d'un unique élément
$\xi \in L(B)$. Notons que $\xi$ possède la propriété universelle suivante :
pour toute $A$-algèbre $C$ et tout $\xi' \in L(C)$, il existe un homomorphisme de
$A$-algèbres $B \to C$ tel que $\xi$ ait pour image $\xi'$ par $L(B) \to L(C)$.

\medskip\noindent
Notons que $B$ est une $A$-algèbre graduée ; les
lemmes \ref{lemma-flat-functor-split} et \ref{lemma-adequate-flat}
permettent donc de décomposer les valeurs de nos foncteurs en $B$ en composantes graduées.
Notons que $\xi \in L(B)^{(1)}$, puisque $(x_1, \ldots, x_n)$ est un élément
de degré un de $\underline{A^{\oplus n}}(B)$. On voit donc que
$\varphi(\xi) \in \underline{M}(B)^{(1)} = M \otimes_A B_1$.
Puisque $B_1$ est engendré par $x_1, \ldots, x_n$ comme $A$-module,
on peut écrire $\varphi(\xi) = \sum m_i \otimes x_i$. Considérons l'application
$A^{\oplus n} \to M$ qui envoie le $i$-ième vecteur de base sur $m_i$.
Par construction, l'application associée
$\underline{A^{\oplus n}} \to \underline{M}$
envoie l'élément $\xi$ sur $\varphi(\xi)$. La propriété
universelle mentionnée ci-dessus entraîne que le diagramme commute.
\end{proof}

\begin{lemma}
\label{lemma-cokernel-into-module}
Soit $A$ un anneau.
Soit $\varphi : F \to \underline{M}$ une application de foncteurs à valeurs dans les modules
sur $\textit{Alg}_A$, avec $F$ adéquat.
Alors $\Coker(\varphi)$ est adéquat.
\end{lemma}

\begin{proof}
D'après le
lemme \ref{lemma-adequate-surjection-from-linear},
on peut supposer que $F = \bigoplus L_i$ est une somme directe de foncteurs
linéairement adéquats. Choisissons des suites exactes
$0 \to L_i \to \underline{A^{\oplus n_i}} \to \underline{A^{\oplus m_i}}$.
Pour chaque $i$, choisissons une application $A^{\oplus n_i} \to M$ comme dans le
lemme \ref{lemma-lift-map}.
Considérons le diagramme
$$
\xymatrix{
0 \ar[r] &
\bigoplus L_i  \ar[r] \ar[d] &
\bigoplus \underline{A^{\oplus n_i}} \ar[r] \ar[ld] &
\bigoplus \underline{A^{\oplus m_i}} \\
& \underline{M}
}
$$
Considérons les $A$-modules
$$
Q =
\Coker(\bigoplus A^{\oplus n_i} \to M \oplus \bigoplus A^{\oplus m_i})
\quad\text{et}\quad
P = \Coker(\bigoplus A^{\oplus n_i} \to \bigoplus A^{\oplus m_i}).
$$
On voit alors que $\Coker(\varphi)$ est isomorphe au
noyau de $\underline{Q} \to \underline{P}$.
\end{proof}

\begin{lemma}
\label{lemma-cokernel-adequate}
\begin{slogan}
Le conoyau d'une application de foncteurs adéquats sur la catégorie des algèbres
sur un anneau est adéquat.
\end{slogan}
Soit $A$ un anneau.
Soit $\varphi : F \to G$ une application de foncteurs adéquats sur $\textit{Alg}_A$.
Alors $\Coker(\varphi)$ est adéquat.
\end{lemma}

\begin{proof}
Choisissons une injection $G \to \underline{M}$.
On a alors une injection $G/F \to \underline{M}/F$. D'après le
lemme \ref{lemma-cokernel-into-module},
$\underline{M}/F$ est adéquat ; on peut donc trouver une injection
$\underline{M}/F \to \underline{N}$.
Par composition, on obtient une injection $G/F \to \underline{N}$. D'après le
lemme \ref{lemma-cokernel-into-module},
le conoyau de l'application induite $G \to \underline{N}$ est adéquat ;
on peut donc trouver une injection $\underline{N}/G \to \underline{K}$.
La suite $0 \to G/F \to \underline{N} \to \underline{K}$ est alors exacte,
ce qui conclut.
\end{proof}

\begin{lemma}
\label{lemma-kernel-adequate}
Soit $A$ un anneau.
Soit $\varphi : F \to G$ une application de foncteurs adéquats sur $\textit{Alg}_A$.
Alors $\Ker(\varphi)$ est adéquat.
\end{lemma}

\begin{proof}
Choisissons une injection $F \to \underline{M}$ et une injection
$G \to \underline{N}$. Notons $F \to \underline{M \oplus N}$
l'application diagonale, de sorte que
$$
\xymatrix{
F \ar[d] \ar[r] & G \ar[d] \\
\underline{M \oplus N} \ar[r] & \underline{N}
}
$$
commute. D'après le
lemme \ref{lemma-cokernel-adequate},
on peut trouver une application de modules $M \oplus N \to K$ telle que
$F$ soit le noyau de $\underline{M \oplus N} \to \underline{K}$.
Alors $\Ker(\varphi)$ est le noyau de
$\underline{M \oplus N} \to \underline{K \oplus N}$.
\end{proof}

\begin{lemma}
\label{lemma-colimit-adequate}
Soit $A$ un anneau.
Une somme directe arbitraire de foncteurs adéquats sur $\textit{Alg}_A$
est adéquate. Une limite inductive de foncteurs adéquats est adéquate.
\end{lemma}

\begin{proof}
L'assertion relative aux sommes directes est immédiate.
Une limite inductive quelconque peut s'écrire comme noyau d'une application entre
sommes directes ; voir
Catégories, lemme \ref{categories-lemma-colimits-coproducts-coequalizers}.
L'assertion résulte donc du
lemme \ref{lemma-kernel-adequate}.
\end{proof}

\begin{lemma}
\label{lemma-flat-linear-functor}
Soit $A$ un anneau.
Soient $F, G$ des foncteurs à valeurs dans les modules sur $\textit{Alg}_A$.
Soit $\varphi : F \to G$ une transformation de foncteurs. Supposons que
\begin{enumerate}
\item $\varphi$ soit additive,
\item pour toute $A$-algèbre $B$, tout $\xi \in F(B)$ et toute unité
$u \in B^*$, on ait $\varphi(u\xi) = u\varphi(\xi)$ dans $G(B)$, et
\item pour tout homomorphisme plat d'anneaux $B \to B'$, on ait
$G(B) \otimes_B B' = G(B')$.
\end{enumerate}
Alors $\varphi$ est un morphisme de foncteurs à valeurs dans les modules.
\end{lemma}

\begin{proof}
Soient $B$ une $A$-algèbre, $\xi \in F(B)$ et $b \in B$. Il faut montrer
que $\varphi(b \xi) = b \varphi(\xi)$. Considérons l'homomorphisme d'anneaux
$$
B \to B' = B[x, y, x^{-1}, y^{-1}]/(x + y - b).
$$
Cet homomorphisme est fidèlement plat, donc $G(B) \subset G(B')$. D'autre
part,
$$
\varphi(b\xi) = \varphi((x + y)\xi) =
\varphi(x\xi) + \varphi(y\xi) = x\varphi(\xi) + y\varphi(\xi)
= (x + y)\varphi(\xi) = b\varphi(\xi)
$$
car $x, y$ sont des unités de $B'$. Cela conclut.
\end{proof}

\begin{lemma}
\label{lemma-extension-adequate-key}
Soit $A$ un anneau.
Soit $0 \to \underline{M} \to G \to L \to 0$ une suite exacte courte
de foncteurs à valeurs dans les modules sur $\textit{Alg}_A$, avec $L$ linéairement adéquat.
Alors $G$ est adéquat.
\end{lemma}

\begin{proof}
Remarquons d'abord que, pour tout homomorphisme plat de $A$-algèbres
$B \to B'$, l'application $G(B) \otimes_B B' \to G(B')$ est un isomorphisme.
Cela vaut en effet pour $\underline{M}$ et $L$, voir le
lemme \ref{lemma-adequate-flat},
et en résulte pour $G$ par le lemme des cinq. En particulier, le
lemme \ref{lemma-flat-functor-split}
montre que $G(B) = \bigoplus_{k \in \mathbf{Z}} G(B)^{(k)}$
pour toute $A$-algèbre graduée $B$.

\medskip\noindent
Choisissons une suite exacte
$0 \to L \to \underline{A^{\oplus n}} \to \underline{A^{\oplus m}}$.
Supposons que l'application $A^{\oplus n} \to A^{\oplus m}$ soit donnée par la
matrice de taille $m \times n$ notée $(a_{ij})$. Considérons la $A$-algèbre graduée
$B = A[x_1, \ldots, x_n]/(\sum a_{ij}x_j)$. L'élément
$(x_1, \ldots, x_n) \in \underline{A^{\oplus n}}(B)$ a une image nulle dans
$\underline{A^{\oplus m}}(B)$ ; il est donc l'image d'un unique élément
$\xi \in L(B)$. Observons que $\xi \in L(B)^{(1)}$. L'application
$$
\Hom_A(B, C) \longrightarrow L(C), \quad
f \longmapsto L(f)(\xi)
$$
définit un isomorphisme de foncteurs. En effet, $f$ est
déterminé par les images $c_i = f(x_i) \in C$, qui doivent
satisfaire aux relations $\sum a_{ij}c_j = 0$. Or $L(C)$ est
l'ensemble des $n$-uplets $(c_1, \ldots, c_n)$ satisfaisant aux relations
$\sum a_{ij} c_j = 0$.

\medskip\noindent
Puisque la valeur de chacun des foncteurs $\underline{M}$, $G$, $L$
en $B$ est la somme directe de ses espaces de poids (d'après le lemme
mentionné ci-dessus), l'exactitude de $0 \to \underline{M} \to G \to L \to 0$ entraîne
celle de la suite $0 \to \underline{M}(B)^{(1)} \to G(B)^{(1)} \to L(B)^{(1)} \to 0$.
On peut donc choisir un élément $\theta \in G(B)^{(1)}$ ayant pour image
$\xi$.

\medskip\noindent
Considérons la $A$-algèbre graduée
$$
C = A[x_1, \ldots, x_n, y_1, \ldots, y_n]/
(\sum a_{ij}x_j, \sum a_{ij}y_j)
$$
Il existe trois homomorphismes de $A$-algèbres graduées $p_1, p_2, m : B \to C$
définis par les règles
$$
p_1(x_i) = x_i, \quad
p_1(x_i) = y_i, \quad
m(x_i) = x_i + y_i.
$$
Nous montrerons que l'élément
$$
\tau = G(m)(\theta) - G(p_1)(\theta) - G(p_2)(\theta) \in G(C)
$$
est nul. Tout d'abord, un calcul direct montre que $\tau$ a une image nulle dans $L(C)$.
Ainsi, $\tau$ est un élément de $\underline{M}(C)$.
De plus, puisque $m$, $p_1$, $p_2$ sont des homomorphismes d'algèbres graduées,
on voit que $\tau \in G(C)^{(1)}$ ; comme $\underline{M} \subset G$,
on en conclut que
$$
\tau \in \underline{M}(C)^{(1)} = M \otimes_A C_1.
$$
On peut écrire de manière unique
$\tau = \underline{M}(p_1)(\tau_1) + \underline{M}(p_2)(\tau_2)$, avec
$\tau_i \in M \otimes_A B_1 = \underline{M}(B)^{(1)}$, car
$C_1 = p_1(B_1) \oplus p_2(B_1)$.
Considérons l'homomorphisme d'anneaux $q_1 : C \to B$ défini par $x_i \mapsto x_i$ et
$y_i \mapsto 0$. Alors
$\underline{M}(q_1)(\tau) =
\underline{M}(q_1)(\underline{M}(p_1)(\tau_1) + \underline{M}(p_2)(\tau_2)) =
\tau_1$.
D'autre part, puisque
$q_1 \circ m = q_1 \circ p_1$, on voit que
$G(q_1)(\tau) = - G(q_1 \circ p_2)(\tau)$. Comme $q_1 \circ p_2$ se factorise en
$B \to A \to B$, on voit que $G(q_1 \circ p_2)(\tau)$ appartient à
$G(B)^{(0)}$ ; voir le
lemme \ref{lemma-flat-functor-split}.
Ainsi $\tau_1 = 0$, car cet élément appartient à
$G(B)^{(0)} \cap \underline{M}(B)^{(1)} \subset
G(B)^{(0)} \cap G(B)^{(1)} = 0$.
De même $\tau_2 = 0$, d'où $\tau = 0$.

\medskip\noindent
Puisque $\theta \in G(B)$, on obtient une transformation de foncteurs
$$
\psi : L(-) = \Hom_A(B, - ) \longrightarrow G(-)
$$
en envoyant $f : B \to C$ sur $G(f)(\theta)$. Puisque $\theta$ est un relèvement de
$\xi$, l'application $\psi$ est un inverse à droite de $G \to L$. En termes de
$\psi$, les assertions démontrées ci-dessus signifient ceci :
$\tau = 0$ signifie que $\psi$ est additive, et
$\theta \in G(B)^{(1)}$ entraîne que, pour toute $A$-algèbre $D$, on a
$\psi(ul) = u\psi(l)$ dans $G(D)$ pour $l \in L(D)$ et toute unité $u \in D^*$.
Cela entraîne que $\psi$ est un morphisme de foncteurs à valeurs dans les modules ; voir le
lemme \ref{lemma-flat-linear-functor}.
Il en résulte clairement que $G \cong \underline{M} \oplus L$, ce qui conclut.
\end{proof}

\begin{remark}
\label{remark-linearly-adequate}
Soit $A$ un anneau.
La démonstration du
lemme \ref{lemma-extension-adequate-key}
montre que toute extension $0 \to \underline{M} \to E \to L \to 0$
de foncteurs à valeurs dans les modules sur $\textit{Alg}_A$,
avec $L$ linéairement adéquat, est scindée. Elle n'utilise que les propriétés suivantes
du foncteur à valeurs dans les modules $F = \underline{M}$ :
\begin{enumerate}
\item $F(B) \otimes_B B' \to F(B')$ est un isomorphisme
pour un homomorphisme plat d'anneaux $B \to B'$, et
\item
$F(C)^{(1)} = F(p_1)(F(B)^{(1)}) \oplus F(p_2)(F(B)^{(1)})$,
où $B = A[x_1, \ldots, x_n]/(\sum a_{ij}x_j)$ et
$C = A[x_1, \ldots, x_n, y_1, \ldots, y_n]/
(\sum a_{ij}x_j, \sum a_{ij}y_j)$.
\end{enumerate}
Ces deux propriétés valent pour tout foncteur adéquat $F$ ; nous omettons les détails.
On voit donc que $L$ est un objet projectif de la catégorie abélienne des
foncteurs adéquats.
\end{remark}

\begin{lemma}
\label{lemma-extension-adequate}
Soit $A$ un anneau.
Soit $0 \to F \to G \to H \to 0$ une suite exacte courte de
foncteurs à valeurs dans les modules sur $\textit{Alg}_A$.
Si $F$ et $H$ sont adéquats, $G$ l'est aussi.
\end{lemma}

\begin{proof}
Choisissons une suite exacte $0 \to F \to \underline{M} \to \underline{N}$.
Si l'on peut montrer que $(\underline{M} \oplus G)/F$ est adéquat, alors
$G$ est le noyau de l'application de foncteurs adéquats
$(\underline{M} \oplus G)/F \to \underline{N}$, donc est
adéquat d'après le
lemme \ref{lemma-kernel-adequate}.
On peut donc supposer que $F = \underline{M}$.

\medskip\noindent
On peut choisir une surjection $L \to H$, où $L$ est une somme directe de
foncteurs linéairement adéquats ; voir le
lemme \ref{lemma-adequate-surjection-from-linear}.
Si l'on peut montrer que le produit fibré $G \times_H L$ est adéquat, alors
$G$ est le conoyau de l'application $\Ker(L \to H) \to G \times_H L$,
donc est adéquat d'après le
lemme \ref{lemma-cokernel-adequate}.
On peut donc supposer que $H = \bigoplus L_i$ est une somme directe de
foncteurs linéairement adéquats. D'après le
lemme \ref{lemma-extension-adequate-key},
chacun des produits fibrés $G \times_H L_i$ est adéquat. D'après le
lemme \ref{lemma-colimit-adequate},
on voit que $\bigoplus G \times_H L_i$ est adéquat.
Alors $G$ est le conoyau de
$$
\bigoplus\nolimits_{i \not = i'} F \longrightarrow
\bigoplus G \times_H L_i
$$
où $\xi$, dans le facteur direct $(i, i')$, a pour image
$(0, \ldots, 0, \xi, 0, \ldots, 0, -\xi, 0, \ldots, 0)$,
avec des composantes non nulles dans les facteurs directs $i$ et $i'$.
Ainsi $G$ est adéquat d'après le
lemme \ref{lemma-cokernel-adequate}.
\end{proof}

\begin{lemma}
\label{lemma-base-change-adequate}
Soit $A \to A'$ un homomorphisme d'anneaux. Si $F$ est un foncteur adéquat sur
$\textit{Alg}_A$, sa restriction $F'$ à
$\textit{Alg}_{A'}$ est également adéquate.
\end{lemma}

\begin{proof}
Choisissons une suite exacte $0 \to F \to \underline{M} \to \underline{N}$.
Alors $F'(B') = F(B') = \Ker(M \otimes_A B' \to N \otimes_A B')$.
Puisque $M \otimes_A B' = M \otimes_A A' \otimes_{A'} B'$, et de même
pour $N$, on voit que $F'$ est le noyau de
$\underline{M \otimes_A A'} \to \underline{N \otimes_A A'}$.
\end{proof}

\begin{lemma}
\label{lemma-pushforward-adequate}
Soit $A \to A'$ un homomorphisme d'anneaux. Si $F'$ est un foncteur adéquat sur
$\textit{Alg}_{A'}$, le foncteur à valeurs dans les modules
$F : B \mapsto F'(A' \otimes_A B)$ sur $\textit{Alg}_A$ est également adéquat.
\end{lemma}

\begin{proof}
Choisissons une suite exacte $0 \to F' \to \underline{M'} \to \underline{N'}$.
Alors
\begin{align*}
F(B) & = F'(A' \otimes_A B) \\
& = \Ker(M' \otimes_{A'} (
A' \otimes_A B) \to N' \otimes_{A'} (A' \otimes_A B)) \\
& = \Ker(M' \otimes_A B \to N' \otimes_A B)
\end{align*}
Ainsi $F$ est le noyau de
$\underline{M} \to \underline{N}$,
où $M = M'$ et $N = N'$ sont considérés comme $A$-modules.
\end{proof}

\begin{lemma}
\label{lemma-adequate-product}
Soit $A = A_1 \times \ldots \times A_n$ un produit d'anneaux.
Se donner un foncteur adéquat sur $A$ revient à se donner une suite
$F_1, \ldots, F_n$ de foncteurs adéquats $F_i$ sur $A_i$.
\end{lemma}

\begin{proof}
Cela résulte de ce qu'une $A$-algèbre $B$ est canoniquement un produit
$B_1 \times \ldots \times B_n$, et qu'il en va de même des $A$-modules.
La correspondance est donnée par $F(B) = \coprod F_i(B_i)$.
Nous omettons les détails.
\end{proof}

\begin{lemma}
\label{lemma-adequate-descent}
Soient $A \to A'$ un homomorphisme d'anneaux et $F$ un foncteur à valeurs dans les modules sur
$\textit{Alg}_A$ tel que
\begin{enumerate}
\item la restriction $F'$ de $F$ à la catégorie des $A'$-algèbres soit
adéquate, et
\item pour toute $A$-algèbre $B$, la suite
$$
0 \to F(B) \to F(B \otimes_A A') \to F(B \otimes_A A' \otimes_A A')
$$
soit exacte.
\end{enumerate}
Alors $F$ est adéquat.
\end{lemma}

\begin{proof}
Les foncteurs $B \to F(B \otimes_A A')$ et
$B \mapsto F(B \otimes_A A' \otimes_A A')$ sont adéquats ; voir les
lemmes \ref{lemma-pushforward-adequate} et
\ref{lemma-base-change-adequate}.
Ainsi $F$, noyau d'une application de foncteurs adéquats, est adéquat ; voir le
lemme \ref{lemma-kernel-adequate}.
\end{proof}






\section{Groupes Ext supérieurs des foncteurs adéquats}
\label{section-higher-ext}

\noindent
Soit $A$ un anneau. Dans le
lemme \ref{lemma-extension-adequate},
nous avons vu que toute extension de foncteurs adéquats dans la catégorie
des foncteurs à valeurs dans les modules sur $\textit{Alg}_A$ est adéquate. Dans cette
section, nous montrons que la même propriété vaut pour les groupes Ext supérieurs.

\begin{lemma}
\label{lemma-adjoint}
Soit $A$ un anneau.
Pour tout foncteur à valeurs dans les modules $F$ sur $\textit{Alg}_A$,
il existe un morphisme $Q(F) \to F$ de foncteurs à valeurs dans les modules sur
$\textit{Alg}_A$ tel que (1) $Q(F)$ soit adéquat et (2) pour tout
foncteur adéquat $G$, l'application $\Hom(G, Q(F)) \to \Hom(G, F)$
soit bijective.
\end{lemma}

\begin{proof}
Choisissons un ensemble $\{L_i\}_{i \in I}$ de foncteurs linéairement adéquats tel que
tout foncteur linéairement adéquat soit isomorphe à l'un des $L_i$.
C'est possible. Supposons que l'on puisse trouver $Q(F) \to F$ satisfaisant à (1) et à la condition
(2)' pour tout $i \in I$, l'application $\Hom(L_i, Q(F)) \to \Hom(L_i, F)$
est bijective. Alors (2) est vérifiée. En effet, en combinant les
lemmes \ref{lemma-adequate-surjection-from-linear} et
\ref{lemma-kernel-adequate},
on voit que tout foncteur adéquat $G$ s'insère dans une suite exacte
$$
K \to L \to G \to 0
$$
où $K$ et $L$ sont des sommes directes de foncteurs linéairement adéquats. Ainsi (2)'
entraîne que
$\Hom(L, Q(F)) \to \Hom(L, F)$
et
$\Hom(K, Q(F)) \to \Hom(K, F)$
sont bijectives, d'où la même assertion pour $G$.

\medskip\noindent
Considérons la catégorie $\mathcal{I}$ dont les objets sont les couples
$(i, \varphi)$, où $i \in I$ et où $\varphi : L_i \to F$ est un morphisme.
Un morphisme $(i, \varphi) \to (i', \varphi')$ est une application
$\psi : L_i \to L_{i'}$ telle que $\varphi' \circ \psi = \varphi$.
Posons
$$
Q(F) = \colim_{(i, \varphi) \in \Ob(\mathcal{I})} L_i
$$
Il existe une application naturelle $Q(F) \to F$ ; d'après le
lemme \ref{lemma-colimit-adequate},
ce foncteur est adéquat et, par construction, il possède la propriété (2)'.
\end{proof}

\begin{lemma}
\label{lemma-enough-injectives}
Soit $A$ un anneau. Notons $\mathcal{P}$ la catégorie des foncteurs à valeurs
dans les modules sur $\textit{Alg}_A$ et $\mathcal{A}$ la catégorie des foncteurs
adéquats sur $\textit{Alg}_A$. Notons $i : \mathcal{A} \to \mathcal{P}$
le foncteur d'inclusion. Notons $Q : \mathcal{P} \to \mathcal{A}$
la construction du lemme \ref{lemma-adjoint}.
Alors
\begin{enumerate}
\item $i$ est pleinement fidèle, exact, et son image est une sous-catégorie de Serre faible,
\item $\mathcal{P}$ possède suffisamment d'objets injectifs,
\item le foncteur $Q$ est adjoint à droite de $i$, donc exact à gauche,
\item $Q$ transforme les objets injectifs en objets injectifs,
\item $\mathcal{A}$ possède suffisamment d'objets injectifs.
\end{enumerate}
\end{lemma}

\begin{proof}
Ce lemme rassemble simplement des faits déjà établis.
L'assertion (1) résulte des définitions, de la caractérisation des
sous-catégories de Serre faibles (voir
Homologie, lemme \ref{homology-lemma-characterize-weak-serre-subcategory})
et des
lemmes \ref{lemma-cokernel-adequate}, \ref{lemma-kernel-adequate}
et \ref{lemma-extension-adequate}.
Rappelons que $\mathcal{P}$ est équivalente à la catégorie
$\textit{PMod}((\textit{Aff}/\Spec(A))_\tau, \mathcal{O})$.
D'où (2), d'après
Injectifs, proposition \ref{injectives-proposition-presheaves-modules}.
L'assertion (3) résulte du
lemme \ref{lemma-adjoint}
et de
Catégories, lemme \ref{categories-lemma-adjoint-exact}.
Les assertions (4) et (5) résultent de
Homologie, lemmes \ref{homology-lemma-adjoint-preserve-injectives} et
\ref{homology-lemma-adjoint-enough-injectives}.
\end{proof}

\noindent
Soit $A$ un anneau. Comme dans
Théorie des déformations formelles, section
\ref{formal-defos-section-tangent-spaces-functors},
étant donnés une $A$-algèbre $B$ et un $B$-module $N$, nous notons $B[N]$
la $R$-algèbre dont le $B$-module sous-jacent est $B \oplus N$ et dont la multiplication
est donnée par $(b, m)(b', m ') = (bb', bm' + b'm)$. Notons que cette construction
est fonctorielle par rapport au couple $(B, N)$, un morphisme $(B, N) \to (B', N')$
étant donné par un homomorphisme de $A$-algèbres $B \to B'$ et une application de $B$-modules
$N \to N'$. En un certain sens, le foncteur $TF$ sur les couples, défini dans le lemme
suivant, est l'espace tangent de $F$.
Ci-dessous, nous ne considérerons que des couples $(B, N)$ tels que
$B[N]$ soit un objet de $\textit{Alg}_A$.

\begin{lemma}
\label{lemma-tangent-functor}
Soit $A$ un anneau. Soit $F$ un foncteur à valeurs dans les modules.
Pour tout $B \in \Ob(\textit{Alg}_A)$ et tout $B$-module $N$,
il existe une décomposition canonique
$$
F(B[N]) = F(B) \oplus TF(B, N)
$$
caractérisée par les propriétés suivantes
\begin{enumerate}
\item $TF(B, N) = \Ker(F(B[N]) \to F(B))$,
\item il existe une structure de $B$-module sur $TF(B, N)$
compatible avec la structure de $B[N]$-module sur $F(B[N])$,
\item $TF$ est un foncteur défini sur la catégorie des couples $(B, N)$,
\item
\label{item-mult-map}
il existe des applications canoniques $N \otimes_B F(B) \to TF(B, N)$
induisant une transformation entre foncteurs définis sur la catégorie
des couples $(B, N)$,
\item $TF(B, 0) = 0$, et l'application $TF(B, N) \to TF(B, N')$ est
nulle lorsque $N \to N'$ est l'application nulle.
\end{enumerate}
\end{lemma}

\begin{proof}
Puisque $B \to B[N] \to B$ est l'identité, on voit que $F(B) \to F(B[N])$
est un facteur direct dont le supplémentaire est $TF(N, B)$, défini en (1).
Cette construction est fonctorielle par rapport au couple $(B, N)$, simplement parce que,
étant donné un morphisme de couples $(B, N) \to (B', N')$, on obtient un diagramme
commutatif
$$
\xymatrix{
B' \ar[r] & B'[N'] \ar[r] & B' \\
B \ar[r] \ar[u] & B[N] \ar[r] \ar[u] & B \ar[u]
}
$$
dans $\textit{Alg}_A$. La structure de $B$-module provient de la structure de $B[N]$-module
et de l'homomorphisme d'anneaux $B \to B[N]$. L'application de (4) est le
composé
$$
N \otimes_B F(B) \longrightarrow
B[N] \otimes_{B[N]} F(B[N]) \longrightarrow F(B[N])
$$
dont l'image est contenue dans $TF(B, N)$. (La première flèche utilise les inclusions
$N \to B[N]$ et $F(B) \to F(B[N])$, et la seconde est l'application
de multiplication.) Si $N = 0$, alors $B = B[N]$,
donc $TF(B, 0) = 0$. Si $N \to N'$ est nulle, elle se factorise en
$N \to 0 \to N'$ ; l'application induite est donc nulle puisque $TF(B, 0) = 0$.
\end{proof}

\noindent
Soit $A$ un anneau. Soit $M$ un $A$-module. Le foncteur à valeurs dans les modules
$\underline{M}$ a pour espace tangent $T\underline{M}$, défini par la règle
$T\underline{M}(B, N) = N \otimes_A M$. En particulier, pour $B$ donné, le
foncteur $N \mapsto T\underline{M}(B, N)$ est additif et exact à droite. Il se trouve
que cela vaut également pour les foncteurs injectifs à valeurs dans les modules.

\begin{lemma}
\label{lemma-tangent-injective}
Soit $A$ un anneau. Soit $I$ un objet injectif de la catégorie
des foncteurs à valeurs dans les modules. Alors, pour tout $B \in \Ob(\textit{Alg}_A)$
et toute suite exacte courte
$0 \to N_1 \to N \to N_2 \to 0$
de $B$-modules, la suite
$$
TI(B, N_1) \to TI(B, N) \to TI(B, N_2) \to 0
$$
est exacte.
\end{lemma}

\begin{proof}
Nous utiliserons les résultats du
lemme \ref{lemma-tangent-functor}
sans le mentionner à nouveau.
Notons $h : \textit{Alg}_A \to \textit{Sets}$ le foncteur donné par
$h(C) = \Mor_A(B[N], C)$. De même pour $h_1$ et $h_2$.
L'application $B[N] \to B[N_2]$ correspondant à la surjection $N \to N_2$
est surjective. Il lui correspond une application $h_2 \to h$ telle que
$h_2(C) \to h(C)$ soit injective pour toutes les $A$-algèbres $C$. D'autre
part, il existe deux applications $p, q : h \to h_1$, correspondant à
l'application nulle $N_1 \to N$ et à l'injection $N_1 \to N$. Notons que
$$
\xymatrix{
h_2 \ar[r] & h \ar@<1ex>[r] \ar@<-1ex>[r] & h_1
}
$$
est un diagramme de noyau double. Notons $\mathcal{O}_h$ le foncteur à valeurs dans les modules
$C \mapsto \bigoplus_{h(C)} C$. De même pour $\mathcal{O}_{h_1}$ et
$\mathcal{O}_{h_2}$. Notons que
$$
\Hom_\mathcal{P}(\mathcal{O}_h, F) = F(B[N])
$$
où $\mathcal{P}$ est la catégorie des foncteurs à valeurs dans les modules sur
$\textit{Alg}_A$. Nous affirmons qu'il existe un diagramme de noyau double
$$
\xymatrix{
\mathcal{O}_{h_2} \ar[r] &
\mathcal{O}_h \ar@<1ex>[r] \ar@<-1ex>[r] &
\mathcal{O}_{h_1}
}
$$
dans $\mathcal{P}$. En effet, supposons que $C \in \Ob(\textit{Alg}_A)$
et que $\xi = \sum_{i = 1, \ldots, n} c_i \cdot f_i$, où $c_i \in C$ et
$f_i : B[N] \to C$, soit un élément de
$\mathcal{O}_h(C)$. Si $p(\xi) = q(\xi)$, alors
on voit que
$$
\sum c_i \cdot f_i \circ z = \sum c_i \cdot f_i \circ y
$$
où $z, y : B[N_1] \to B[N]$ sont les applications $z : (b, m_1) \mapsto (b, 0)$
et $y : (b, m_1) \mapsto (b, m_1)$. Cela signifie que, pour tout $i$,
il existe un $j$ tel que $f_j \circ z = f_i \circ y$. Il en résulte clairement
que $f_i(N_1) = 0$, c'est-à-dire que $f_i$ se factorise par une unique application
$\overline{f}_i : B[N_2] \to C$. Ainsi $\xi$ est l'image de
$\overline{\xi} = \sum c_i \cdot \overline{f}_i$.
Puisque $I$ est injectif, il transforme ce diagramme de noyau double
en un diagramme de conoyau double
$$
\xymatrix{
I(B[N_1]) \ar@<1ex>[r] \ar@<-1ex>[r] &
I(B[N]) \ar[r] &
I(B[N_2])
}
$$
Ce diagramme est compatible avec les décompositions en somme directe
$I(B[N]) = I(B) \oplus TI(B, N)$ et $I(B[N_i]) = I(B) \oplus TI(B, N_i)$.
L'application nulle $N \to N_1$ induit l'application nulle $TI(B, N) \to TI(B, N_1)$.
On voit ainsi que la propriété de conoyau double
ci-dessus signifie que l'on a une suite exacte
$TI(B, N_1) \to TI(B, N) \to TI(B, N_2) \to 0$,
comme voulu.
\end{proof}

\begin{lemma}
\label{lemma-exactness-implies}
Soit $A$ un anneau. Soit $F$ un foncteur à valeurs dans les modules
tel que, pour tout $B \in \Ob(\textit{Alg}_A)$, le
foncteur $TF(B, -)$ sur les $B$-modules transforme une suite exacte courte
de $B$-modules en une suite exacte à droite. Alors
\begin{enumerate}
\item $TF(B, N_1 \oplus N_2) = TF(B, N_1) \oplus TF(B, N_2)$,
\item il existe une seconde structure fonctorielle de $B$-module sur $TF(B, N)$,
définie en posant $x \cdot b = TF(B, b\cdot 1_N)(x)$ pour $x \in TF(B, N)$
et $b \in B$,
\item
\label{item-mult-map-linear}
l'application canonique $N \otimes_B F(B) \to TF(B, N)$ du
lemme \ref{lemma-tangent-functor}
est $B$-linéaire également pour la seconde structure de $B$-module,
\item
\label{item-tangent-right-exact}
étant donné un $B$-module de présentation finie $N$, il existe un
isomorphisme canonique $TF(B, B) \otimes_B N \to TF(B, N)$, où le produit
tensoriel utilise la seconde structure de $B$-module sur $TF(B, B)$.
\end{enumerate}
\end{lemma}

\begin{proof}
Nous utiliserons les résultats du
lemme \ref{lemma-tangent-functor}
sans le mentionner à nouveau.
Les applications $N_1 \to N_1 \oplus N_2$ et $N_2 \to N_1 \oplus N_2$ donnent
une application $TF(B, N_1) \oplus TF(B, N_2) \to TF(B, N_1 \oplus N_2)$,
qui est injective puisque les applications $N_1 \oplus N_2 \to N_1$ et
$N_1 \oplus N_2 \to N_2$ induisent un inverse.
Puisque $TF$ est exact à droite, on voit que
$TF(B, N_1) \to TF(B, N_1 \oplus N_2) \to TF(B, N_2) \to 0$ est exacte.
Ainsi $TF(B, N_1) \oplus TF(B, N_2) \to TF(B, N_1 \oplus N_2)$ est un
isomorphisme. Cela démontre (1).

\medskip\noindent
Pour (2), il suffit de montrer que
$x \cdot (b_1 + b_2) = x \cdot b_1 + x \cdot b_2$.
(L'associativité et l'additivité sont claires.) Pour cela, considérons
$$
N \xrightarrow{(b_1, b_2)} N \oplus N \xrightarrow{+} N
$$
et appliquons $TF(B, -)$.

\medskip\noindent
L'assertion (3) résulte immédiatement de ce que
$N \otimes_B F(B) \to TF(B, N)$ est fonctorielle par rapport au couple $(B, N)$.

\medskip\noindent
Supposons que $N$ soit un $B$-module de présentation finie. Choisissons une présentation
$B^{\oplus m} \to B^{\oplus n} \to N \to 0$. Elle donne une suite
exacte
$$
TF(B, B^{\oplus m}) \to TF(B, B^{\oplus n}) \to TF(B, N) \to 0
$$
par exactitude à droite de $TF(B, -)$. D'après (1), on peut écrire
$TF(B, B^{\oplus m}) = TF(B, B)^{\oplus m}$ et
$TF(B, B^{\oplus n}) = TF(B, B)^{\oplus n}$. Supposons ensuite que
$B^{\oplus m} \to B^{\oplus n}$ soit donnée par la matrice $T = (b_{ij})$.
Alors l'application induite $TF(B, B)^{\oplus m} \to TF(B, B)^{\oplus n}$
est donnée par la matrice de coefficients $TF(B, b_{ij} \cdot 1_B)$.
Ce fait, joint à l'exactitude à droite de $\otimes$, démontre (4).
\end{proof}

\begin{example}
\label{example-module-structure-different}
Soit $F$ un foncteur à valeurs dans les modules comme dans le
lemme \ref{lemma-exactness-implies}.
Les deux structures de module sur
$TF(B, N)$ ne coïncident pas toujours. En voici un exemple. Supposons que $A = \mathbf{F}_p$,
où $p$ est un nombre premier. Posons $F(B) = B$, mais avec la structure de $B$-module
donnée par $b \cdot x = b^px$. Alors $TF(B, N) = N$, avec la structure de $B$-module
donnée par $b \cdot x = b^px$ pour $x \in N$. Cependant, la seconde structure de $B$-module
est donnée par $x \cdot b = bx$. Notons que, dans ce cas, l'application canonique
$N \otimes_B F(B) \to TF(B, N)$ est nulle, car élever un élément
$n \in B[N]$ à la puissance $p$ donne zéro.
\end{example}

\noindent
Dans le lemme suivant, nous utiliserons fréquemment l'observation suivante :
si $0 \to F \to G \to H \to 0$ est une suite exacte de foncteurs à valeurs dans les modules
sur $\textit{Alg}_A$, alors, pour tout couple $(B, N)$, la
suite $0 \to TF(B, N) \to TG(B, N) \to TH(B, N) \to 0$ est exacte.
Cela résulte de ce que $0 \to F(B[N]) \to G(B[N]) \to H(B[N]) \to 0$
est exacte.

\begin{lemma}
\label{lemma-exactness-permanence}
Soit $A$ un anneau. Pour un foncteur à valeurs dans les modules $F$ sur $\textit{Alg}_A$,
disons que $(*)$ est vérifiée si, pour tout $B \in \Ob(\textit{Alg}_A)$, le
foncteur $TF(B, -)$ sur les $B$-modules transforme une suite exacte courte
de $B$-modules en une suite exacte à droite. Soit
$0 \to F \to G \to H \to 0$ une suite exacte courte de
foncteurs à valeurs dans les modules sur $\textit{Alg}_A$.
\begin{enumerate}
\item Si $(*)$ est vérifiée pour $F, G$, alors $(*)$ est vérifiée pour $H$.
\item Si $(*)$ est vérifiée pour $F, H$, alors $(*)$ est vérifiée pour $G$.
\item Si $H' \to H$ est un morphisme de foncteurs à valeurs dans les modules sur $\textit{Alg}_A$
et si $(*)$ est vérifiée pour $F$, $G$, $H$ et $H'$, alors $(*)$ est vérifiée pour
$G \times_H H'$.
\end{enumerate}
\end{lemma}

\begin{proof}
Soit $B$ donné. Soit $0 \to N_1 \to N_2 \to N_3 \to 0$ une suite exacte
courte de $B$-modules. L'assertion (1) résulte d'une chasse dans le
diagramme
$$
\xymatrix{
0 \ar[r] &
TF(B, N_1) \ar[r] \ar[d] &
TG(B, N_1) \ar[r] \ar[d] &
TH(B, N_1) \ar[r] \ar[d] & 0 \\
0 \ar[r] &
TF(B, N_2) \ar[r] \ar[d] &
TG(B, N_2) \ar[r] \ar[d] &
TH(B, N_2) \ar[r] \ar[d] & 0 \\
0 \ar[r] &
TF(B, N_3) \ar[r] \ar[d] &
TG(B, N_3) \ar[r] \ar[d] &
TH(B, N_3) \ar[r] & 0 \\
& 0 & 0
}
$$
dont les lignes horizontales et les colonnes faisant intervenir $TF$ et $TG$ sont exactes.
Pour démontrer (2), on effectue une chasse dans le diagramme
$$
\xymatrix{
0 \ar[r] &
TF(B, N_1) \ar[r] \ar[d] &
TG(B, N_1) \ar[r] \ar[d] &
TH(B, N_1) \ar[r] \ar[d] & 0 \\
0 \ar[r] &
TF(B, N_2) \ar[r] \ar[d] &
TG(B, N_2) \ar[r] \ar[d] &
TH(B, N_2) \ar[r] \ar[d] & 0 \\
0 \ar[r] &
TF(B, N_3) \ar[r] \ar[d] &
TG(B, N_3) \ar[r] &
TH(B, N_3) \ar[r] \ar[d] & 0 \\
& 0 & & 0
}
$$
dont les lignes horizontales et les colonnes faisant intervenir $TF$ et $TH$ sont exactes.
L'assertion (3) résulte de (2), puisque $G \times_H H'$ s'insère dans la suite exacte
$0 \to F \to G \times_H H' \to H' \to 0$.
\end{proof}

\noindent
L'essentiel du travail de cette section a été accompli pour démontrer le
résultat d'annulation fondamental suivant.

\begin{lemma}
\label{lemma-ext-group-zero-key}
Soit $A$ un anneau. Soient $M$, $P$ des $A$-modules, avec $P$ de présentation
finie. Alors
$\Ext^i_\mathcal{P}(\underline{P}, \underline{M}) = 0$
pour $i > 0$, où $\mathcal{P}$ est la catégorie des foncteurs à valeurs dans les modules
sur $\textit{Alg}_A$.
\end{lemma}

\begin{proof}
Choisissons une résolution injective $\underline{M} \to I^\bullet$ dans
$\mathcal{P}$ ; voir le
lemme \ref{lemma-enough-injectives}.
D'après
Catégories dérivées, lemme \ref{derived-lemma-compute-ext-resolutions},
tout élément de $\Ext^i_\mathcal{P}(\underline{P}, \underline{M})$
provient d'un morphisme $\varphi : \underline{P} \to I^i$ tel que
$d^i \circ \varphi = 0$. Nous montrerons que
l'extension de Yoneda
$$
E : 0 \to \underline{M} \to I^0 \to \ldots \to
I^{i - 1} \times_{\Ker(d^i)} \underline{P} \to \underline{P} \to 0
$$
de $\underline{P}$ par $\underline{M}$
associée à $\varphi$ est triviale, ce qui démontrera le lemme d'après
Catégories dérivées, lemme \ref{derived-lemma-yoneda-extension}.

\medskip\noindent
Pour un foncteur à valeurs dans les modules $F$ sur $\textit{Alg}_A$,
disons que $(*)$ est vérifiée si, pour tout $B \in \Ob(\textit{Alg}_A)$, le
foncteur $TF(B, -)$ sur les $B$-modules transforme une suite exacte courte
de $B$-modules en une suite exacte à droite.
Rappelons que les foncteurs à valeurs dans les modules $\underline{M}, I^n, \underline{P}$
possèdent chacun la propriété $(*)$ ; voir le
lemme \ref{lemma-tangent-injective}
et les remarques qui le précèdent.
En décomposant $0 \to \underline{M} \to I^\bullet$ en suites exactes
courtes, on trouve que chacun des foncteurs
$\Im(d^{n - 1}) = \Ker(d^n) \subset I^n$ possède la propriété $(*)$ d'après le
lemme \ref{lemma-exactness-permanence},
et que $I^{i - 1} \times_{\Ker(d^i)} \underline{P}$ possède également la propriété
$(*)$.

\medskip\noindent
On peut donc supposer que l'extension de Yoneda est donnée sous la forme
$$
E : 0 \to \underline{M} \to F_{i - 1} \to \ldots \to
F_0 \to \underline{P} \to 0
$$
où chacun des foncteurs à valeurs dans les modules $F_j$ possède la propriété $(*)$.
Posons $G_j(B) = TF_j(B, B)$, considéré comme $B$-module au moyen de la {\it seconde}
structure de $B$-module définie dans le
lemme \ref{lemma-exactness-implies}.
Puisque $TF_j$ est un foncteur sur les couples, on voit que $G_j$ est un foncteur
à valeurs dans les modules sur $\textit{Alg}_A$. De plus, puisque $E$ est une suite exacte,
la suite $G_{j + 1} \to G_j \to G_{j - 1}$ est exacte (voir la remarque
précédant le
lemme \ref{lemma-exactness-permanence}).
Observons que $T\underline{M}(B, B) = M \otimes_A B = \underline{M}(B)$
et que les deux structures de $B$-module y coïncident.
On obtient ainsi une extension de Yoneda
$$
E' :  0 \to \underline{M} \to G_{i - 1} \to \ldots \to
G_0 \to \underline{P} \to 0
$$
De plus, les applications canoniques
$$
F_j(B) = B \otimes_B F_j(B) \longrightarrow TF_j(B, B) = G_j(B)
$$
du
lemme \ref{lemma-tangent-functor} (\ref{item-mult-map})
sont $B$-linéaires d'après le
lemme \ref{lemma-exactness-implies} (\ref{item-mult-map-linear})
et fonctorielles en $B$. On a donc une application
$$
\xymatrix{
0 \ar[r] &
\underline{M} \ar[r] \ar[d]^1 &
F_{i - 1} \ar[r] \ar[d] &
\ldots \ar[r] &
F_0 \ar[r] \ar[d] &
\underline{P} \ar[r] \ar[d]^1 & 0 \\
0 \ar[r] &
\underline{M} \ar[r] &
G_{i - 1} \ar[r] &
\ldots \ar[r] &
G_0 \ar[r] &
\underline{P} \ar[r] & 0
}
$$
d'extensions de Yoneda. En particulier, on voit que $E$ et $E'$ ont la
même classe dans $\Ext^i_\mathcal{P}(\underline{P}, \underline{M})$,
d'après le lemme sur les Ext de Yoneda mentionné ci-dessus. Enfin, soit $N$ un
$A$-module de présentation finie. On voit alors que
$$
0 \to T\underline{M}(A, N) \to TF_{i - 1}(A, N) \to \ldots \to
TF_0(A, N) \to T\underline{P}(A, N) \to 0
$$
est exacte. D'après le
lemme \ref{lemma-exactness-implies} (\ref{item-tangent-right-exact}),
avec $B = A$, cela se traduit par l'exactitude de la suite de
$A$-modules
$$
0 \to M \otimes_A N \to G_{i - 1}(A) \otimes_A N \to \ldots \to
G_0(A) \otimes_A N \to P \otimes_A N \to 0
$$
Ainsi la suite de $A$-modules
$0 \to M \to G_{i - 1}(A) \to \ldots \to G_0(A) \to P \to 0$
est universellement exacte, au sens où elle reste exacte après tensorisation
par tout $A$-module de présentation finie $N$. Posons
$K = \Ker(G_0(A) \to P)$, de sorte que l'on a des suites exactes
$$
0 \to K \to G_0(A) \to P \to 0
\quad\text{et}\quad
G_2(A) \to G_1(A) \to K \to 0
$$
En tensorisant la seconde suite par $N$, on obtient
$K \otimes_A N = \Coker(G_2(A) \otimes_A N \to G_1(A) \otimes_A N)$.
L'exactitude de $G_2(A) \otimes_A N \to G_1(A) \otimes_A N \to G_0(A) \otimes_A N$
entraîne alors que $K \otimes_A N \to G_0(A) \otimes_A N$ est injective.
D'après
Algèbre, théorème \ref{algebra-theorem-universally-exact-criteria},
cela signifie que l'extension de $A$-modules $0 \to K \to G_0(A) \to P \to 0$
est exacte ; puisque $P$ est supposé de présentation finie, cela signifie
que la suite est scindée ; voir
Algèbre, lemme \ref{algebra-lemma-universally-exact-split}.
Tout scindage $P \to G_0(A)$ définit une application $\underline{P} \to G_0$
qui scinde la surjection $G_0 \to \underline{P}$. Ainsi
l'extension de Yoneda $E'$ est équivalente à l'extension de Yoneda triviale,
ce qui conclut.
\end{proof}

\begin{lemma}
\label{lemma-ext-group-zero}
Soit $A$ un anneau. Soit $M$ un $A$-module. Soit $L$ un foncteur linéairement
adéquat sur $\textit{Alg}_A$. Alors
$\Ext^i_\mathcal{P}(L, \underline{M}) = 0$
pour $i > 0$, où $\mathcal{P}$ est la catégorie des foncteurs à valeurs dans les modules
sur $\textit{Alg}_A$.
\end{lemma}

\begin{proof}
Puisque $L$ est linéairement adéquat, il existe une suite exacte
$$
0 \to L \to \underline{A^{\oplus m}} \to \underline{A^{\oplus n}} \to
\underline{P} \to 0
$$
Ici $P = \Coker(A^{\oplus m} \to A^{\oplus n})$ est le conoyau
de l'application de $A$-modules libres de rang fini donnée par la définition
des foncteurs linéairement adéquats. D'après le
lemme \ref{lemma-ext-group-zero-key},
on a l'annulation de
$\Ext^i_\mathcal{P}(\underline{P}, \underline{M})$
et de
$\Ext^i_\mathcal{P}(\underline{A}, \underline{M})$
pour $i > 0$.
Posons $K = \Ker(\underline{A^{\oplus n}} \to \underline{P})$.
La suite exacte longue des groupes Ext associée à la suite exacte
$0 \to K \to \underline{A^{\oplus n}} \to \underline{P} \to 0$
permet de conclure que
$\Ext^i_\mathcal{P}(K, \underline{M}) = 0$ pour $i > 0$.
En répétant l'argument avec la suite
$0 \to L \to \underline{A^{\oplus m}} \to K \to 0$,
on conclut.
\end{proof}

\begin{lemma}
\label{lemma-RQ-zero}
Avec les notations du
lemme \ref{lemma-enough-injectives},
on a $R^pQ(F) = 0$ pour tout $p > 0$ et tout foncteur adéquat $F$.
\end{lemma}

\begin{proof}
Choisissons une suite exacte $0 \to F \to \underline{M^0} \to \underline{M^1}$.
Posons $M^2 = \Coker(M^0 \to M^1)$, de sorte que
$0 \to F \to \underline{M^0} \to \underline{M^1}
\to \underline{M^2} \to 0$ est une résolution. D'après
Catégories dérivées, lemme \ref{derived-lemma-two-ss-complex-functor},
on obtient une suite spectrale
$$
R^pQ(\underline{M^q}) \Rightarrow R^{p + q}Q(F)
$$
Puisque $Q(\underline{M^q}) = \underline{M^q}$, il suffit de démontrer
$R^pQ(\underline{M}) = 0$, $p > 0$, pour tout $A$-module $M$.

\medskip\noindent
Choisissons une résolution injective $\underline{M} \to I^\bullet$ dans
la catégorie $\mathcal{P}$. Supposons que $R^iQ(\underline{M})$ soit non nul.
Alors $\Ker(Q(I^i) \to Q(I^{i + 1}))$ contient strictement
l'image de $Q(I^{i - 1}) \to Q(I^i)$. Ainsi, d'après le
lemme \ref{lemma-adequate-surjection-from-linear},
il existe un foncteur linéairement adéquat $L$ et une application
$\varphi : L \to Q(I^i)$ à valeurs dans le noyau de $Q(I^i) \to Q(I^{i + 1})$
qui ne se factorise pas par l'image de $Q(I^{i - 1}) \to Q(I^i)$.
Puisque $Q$ est adjoint à gauche du foncteur d'inclusion, l'application
$\varphi$ correspond à une application $\varphi' : L \to I^i$ ayant les mêmes propriétés.
Ainsi $\varphi'$ donne un élément non nul de
$\Ext^i_\mathcal{P}(L, \underline{M})$, en contradiction avec le
lemme \ref{lemma-ext-group-zero}.
\end{proof}














\section{Modules adéquats}
\label{section-adequate}

\noindent
Dans
Descente, section \ref{descent-section-quasi-coherent-sheaves},
nous avons vu que les modules quasi-cohérents sur un schéma $S$
s'identifient aux modules quasi-cohérents sur n'importe lequel des grands
sites $(\Sch/S)_\tau$ associés à $S$. Nous avons vu que cette
identification présente deux difficultés :
\begin{enumerate}
\item $\QCoh(\mathcal{O}_S) \to
\textit{Mod}((\Sch/S)_\tau, \mathcal{O})$,
$\mathcal{F} \mapsto \mathcal{F}^a$ n'est pas exact en général
(Descente, lemme \ref{descent-lemma-equivalence-quasi-coherent-limits}), et
\item étant donné un morphisme quasi-compact et quasi-séparé $f : X \to S$,
le foncteur $f_*$ ne préserve pas en général les faisceaux quasi-cohérents sur les
grands sites (Descente, proposition
\ref{descent-proposition-equivalence-quasi-coherent-functorial}).
\end{enumerate}
L'assertion (1) signifie que l'on ne peut pas définir une sous-catégorie triangulée
de $D(\mathcal{O})$ constituée des complexes dont les faisceaux de cohomologie
sont quasi-cohérents. L'assertion (2) signifie que $Rf_*\mathcal{F}$ n'est pas un
complexe à faisceaux de cohomologie quasi-cohérents, même lorsque $\mathcal{F}$
est quasi-cohérent et que $f$ est quasi-compact et quasi-séparé.
De plus, les exemples donnés dans les démonstrations de
Descente, lemme
\ref{descent-lemma-equivalence-quasi-coherent-limits},
et de
Descente, proposition
\ref{descent-proposition-equivalence-quasi-coherent-functorial},
ne sont pas de nature pathologique.

\medskip\noindent
Dans cette section, nous étudions une catégorie légèrement plus grande
de $\mathcal{O}$-modules sur $(\Sch/S)_\tau$, qui contient les
modules quasi-cohérents, est abélienne et est préservée par $f_*$ lorsque
$f$ est quasi-compact et quasi-séparé.
Pour cela, supposons que $S$ soit un schéma. Soit $\mathcal{F}$ un préfaisceau
de $\mathcal{O}$-modules sur $(\Sch/S)_\tau$.
Pour tout objet affine $U = \Spec(A)$ de $(\Sch/S)_\tau$,
on peut restreindre $\mathcal{F}$ à $(\textit{Aff}/U)_\tau$ pour obtenir
un préfaisceau de $\mathcal{O}$-modules sur ce site. Le foncteur correspondant
à valeurs dans les modules, voir la
section \ref{section-quasi-coherent},
sera noté
$$
F = F_{\mathcal{F}, A} :
\textit{Alg}_A \longrightarrow \textit{Ab},
\quad
B \longmapsto \mathcal{F}(\Spec(B))
$$
La correspondance $\mathcal{F} \mapsto F_{\mathcal{F}, A}$ est un foncteur exact
de catégories abéliennes.

\begin{definition}
\label{definition-adequate}
Un faisceau de $\mathcal{O}$-modules $\mathcal{F}$ sur $(\Sch/S)_\tau$ est
{\it adéquat} s'il existe un $\tau$-recouvrement
$\{\Spec(A_i) \to S\}_{i \in I}$ tel que $F_{\mathcal{F}, A_i}$ soit
adéquat pour tout $i \in I$.
\end{definition}

\noindent
Nous verrons ci-dessous que la catégorie des $\mathcal{O}$-modules adéquats
est indépendante de la topologie $\tau$ choisie.

\begin{lemma}
\label{lemma-adequate-local}
Soit $S$ un schéma. Soit $\mathcal{F}$ un $\mathcal{O}$-module adéquat sur
$(\Sch/S)_\tau$. Pour tout schéma affine $\Spec(A)$ au-dessus de $S$,
le foncteur $F_{\mathcal{F}, A}$ est adéquat.
\end{lemma}

\begin{proof}
Soit $\{\Spec(A_i) \to S\}_{i \in I}$ un $\tau$-recouvrement
tel que $F_{\mathcal{F}, A_i}$ soit adéquat pour tout $i \in I$.
On peut trouver un $\tau$-recouvrement affine standard
$\{\Spec(A'_j) \to \Spec(A)\}_{j = 1, \ldots, m}$
tel que $\Spec(A'_j) \to \Spec(A) \to S$ se factorise
par $\Spec(A_{i(j)})$ pour un certain $i(j) \in I$. On voit alors que
$F_{\mathcal{F}, A'_j}$ est la restriction de
$F_{\mathcal{F}, A_{i(j)}}$ à la catégorie des $A'_j$-algèbres.
Ainsi $F_{\mathcal{F}, A'_j}$ est adéquat d'après le
lemme \ref{lemma-base-change-adequate}.
D'après le
lemme \ref{lemma-adequate-product},
la suite
$F_{\mathcal{F}, A'_j}$ correspond à un foncteur « produit » adéquat
$F'$ sur $A' = A'_1 \times \ldots \times A'_m$. Comme $\mathcal{F}$ est un
faisceau (pour la topologie de Zariski), ce foncteur produit $F'$ est égal
à $F_{\mathcal{F}, A'}$, c'est-à-dire à la restriction de $F$ aux $A'$-algèbres.
Enfin, $\{\Spec(A') \to \Spec(A)\}$ est un $\tau$-recouvrement.
Il résulte du
lemme \ref{lemma-adequate-descent}
que $F_{\mathcal{F}, A}$ est adéquat.
\end{proof}

\begin{lemma}
\label{lemma-adequate-affine}
Soit $S = \Spec(A)$ un schéma affine. La catégorie des
$\mathcal{O}$-modules adéquats sur $(\Sch/S)_\tau$ est équivalente à la
catégorie des foncteurs adéquats à valeurs dans les modules sur $\textit{Alg}_A$.
\end{lemma}

\begin{proof}
Étant donné un module adéquat $\mathcal{F}$, le foncteur $F_{\mathcal{F}, A}$
est adéquat d'après le lemme \ref{lemma-adequate-local}.
Étant donné un foncteur adéquat $F$, choisissons une suite exacte
$0 \to F \to \underline{M} \to \underline{N}$ et considérons
le $\mathcal{O}$-module $\mathcal{F} = \Ker(M^a \to N^a)$, où
$M^a$ désigne le $\mathcal{O}$-module quasi-cohérent sur
$(\Sch/S)_\tau$ associé au faisceau quasi-cohérent
$\widetilde{M}$ sur $S$. Notons que $F = F_{\mathcal{F}, A}$ ; en particulier,
le module $\mathcal{F}$ est adéquat par définition.
Nous omettons la démonstration du fait que ces constructions définissent des équivalences
de catégories inverses l'une de l'autre.
\end{proof}

\begin{lemma}
\label{lemma-pullback-adequate}
Soit $f : T \to S$ un morphisme de schémas.
L'image inverse $f^*\mathcal{F}$ d'un $\mathcal{O}$-module adéquat
$\mathcal{F}$ sur $(\Sch/S)_\tau$ est un
$\mathcal{O}$-module adéquat sur $(\Sch/T)_\tau$.
\end{lemma}

\begin{proof}
Le foncteur d'image inverse
$$
\begin{aligned}
f^* & : \textit{Mod}((\Sch/S)_\tau, \mathcal{O}) \\
& \longrightarrow \textit{Mod}((\Sch/T)_\tau, \mathcal{O})
\end{aligned}
$$
est donné par restriction, c'est-à-dire que $f^*\mathcal{F}(V) = \mathcal{F}(V)$
pour tout schéma $V$ au-dessus de $T$. Ce lemme résulte donc immédiatement du
lemme \ref{lemma-adequate-local}
et de la définition.
\end{proof}

\noindent
Voici une caractérisation de la catégorie des $\mathcal{O}$-modules adéquats.
Pour en comprendre la portée, considérons une application $\mathcal{G} \to \mathcal{H}$
de $\mathcal{O}_S$-modules quasi-cohérents sur un schéma $S$.
Le conoyau de l'application associée $\mathcal{G}^a \to \mathcal{H}^a$
de $\mathcal{O}$-modules est quasi-cohérent, car il est égal à
$(\mathcal{H}/\mathcal{G})^a$. Mais le noyau de
$\mathcal{G}^a \to \mathcal{H}^a$ n'est pas quasi-cohérent
en général. Il est toutefois adéquat.

\begin{lemma}
\label{lemma-adequate-characterize}
Soit $S$ un schéma. Soit $\mathcal{F}$ un $\mathcal{O}$-module sur
$(\Sch/S)_\tau$. Les conditions suivantes sont équivalentes
\begin{enumerate}
\item $\mathcal{F}$ est adéquat,
\item il existe un recouvrement ouvert affine $S = \bigcup S_i$ et des
applications de $\mathcal{O}_{S_i}$-modules quasi-cohérents
$\mathcal{G}_i \to \mathcal{H}_i$
telles que $\mathcal{F}|_{(\Sch/S_i)_\tau}$ soit le
noyau de $\mathcal{G}_i^a \to \mathcal{H}_i^a$
\item il existe un $\tau$-recouvrement $\{S_i \to S\}_{i \in I}$ et des
applications de $\mathcal{O}_{S_i}$-modules quasi-cohérents
$\mathcal{G}_i \to \mathcal{H}_i$
telles que $\mathcal{F}|_{(\Sch/S_i)_\tau}$ soit le
noyau de $\mathcal{G}_i^a \to \mathcal{H}_i^a$,
\item il existe un $\tau$-recouvrement $\{f_i : S_i \to S\}_{i \in I}$
tel que chaque $f_i^*\mathcal{F}$ soit adéquat,
\item pour tout schéma affine $U$ au-dessus de $S$, la restriction
$\mathcal{F}|_{(\Sch/U)_\tau}$ est le noyau
d'une application $\mathcal{G}^a \to \mathcal{H}^a$ de
$\mathcal{O}_U$-modules quasi-cohérents.
\end{enumerate}
\end{lemma}

\begin{proof}
Soit $U = \Spec(A)$ un schéma affine au-dessus de $S$.
Posons $F = F_{\mathcal{F}, A}$. Par définition, le foncteur
$F$ est adéquat si et seulement s'il existe une application de $A$-modules
$M \to N$ telle que $F = \Ker(\underline{M} \to \underline{N})$.
En combinant cela avec les
lemmes \ref{lemma-adequate-local} et
\ref{lemma-adequate-affine},
on voit que (1) et (5) sont équivalentes.

\medskip\noindent
Il est clair que (5) entraîne (2) et que (2) entraîne (3).
Si (3) est vérifiée, on peut raffiner le recouvrement
$\{S_i \to S\}$ de sorte que chaque $S_i = \Spec(A_i)$ soit affine.
Les remarques préliminaires de la démonstration montrent alors que
$F_{\mathcal{F}, A_i}$ est adéquat. Ainsi $\mathcal{F}$
est adéquat par définition. Donc (3) entraîne (1).

\medskip\noindent
Enfin, (4) est équivalente à (1) : on utilise le
lemme \ref{lemma-pullback-adequate}
dans un sens, et le fait qu'un
composé de $\tau$-recouvrements est un $\tau$-recouvrement dans l'autre.
\end{proof}

\noindent
Tout comme pour les faisceaux quasi-cohérents, la catégorie des
modules adéquats est indépendante de la topologie.

\begin{lemma}
\label{lemma-adequate-fpqc}
Soit $\mathcal{F}$ un $\mathcal{O}$-module adéquat sur
$(\Sch/S)_\tau$. Pour tout morphisme plat surjectif
$\Spec(B) \to \Spec(A)$ de schémas affines au-dessus de $S$,
le complexe de {\v C}ech augmenté
$$
0 \to \mathcal{F}(\Spec(A)) \to
\mathcal{F}(\Spec(B)) \to
\mathcal{F}(\Spec(B \otimes_A B)) \to \ldots
$$
est exact. En particulier, $\mathcal{F}$ satisfait à la condition de faisceau
pour les recouvrements fpqc, et c'est un faisceau de $\mathcal{O}$-modules
sur $(\Sch/S)_{fppf}$.
\end{lemma}

\begin{proof}
Avec $A \to B$ comme dans le lemme, posons $F = F_{\mathcal{F}, A}$. Ce foncteur
est adéquat d'après le
lemme \ref{lemma-adequate-local}.
D'après le
lemme \ref{lemma-adequate-flat},
puisque $A \to B$, $A \to B \otimes_A B$, etc., sont plats, on voit que
$F(B) = F(A) \otimes_A B$,
$F(B \otimes_A B) = F(A) \otimes_A B \otimes_A B$, etc.
L'exactitude résulte de
Descente, lemme \ref{descent-lemma-ff-exact}.

\medskip\noindent
Ainsi $\mathcal{F}$ satisfait à la condition de faisceau pour les
$\tau$-recouvrements (en particulier les recouvrements de Zariski) et pour tout
recouvrement fidèlement plat d'un schéma affine par un schéma affine. En raisonnant
comme dans les démonstrations de Descente, lemme \ref{descent-lemma-standard-fpqc-covering},
et de
Descente, proposition \ref{descent-proposition-fpqc-descent-quasi-coherent},
on conclut que $\mathcal{F}$ satisfait à la condition de faisceau pour tous les
recouvrements fpqc (constitués d'objets de $(\Sch/S)_\tau$).
Nous omettons les détails.
\end{proof}

\noindent
Le lemme \ref{lemma-adequate-fpqc} montre en particulier que,
pour tout couple de topologies $\tau, \tau'$, les collections
des modules adéquats pour la topologie $\tau$ et pour la topologie $\tau'$
sont identiques (comme préfaisceaux de modules sur la catégorie sous-jacente $\Sch/S$).

\begin{definition}
\label{definition-category-adequate-modules}
Soit $S$ un schéma. La catégorie des $\mathcal{O}$-modules adéquats sur
$(\Sch/S)_\tau$ est notée {\it $\textit{Adeq}(\mathcal{O})$} ou
{\it $\textit{Adeq}((\Sch/S)_\tau, \mathcal{O})$}. Si l'on veut considérer
la catégorie abélienne des modules adéquats sans choisir de
topologie, on écrit simplement {\it $\textit{Adeq}(S)$}.
\end{definition}

\begin{lemma}
\label{lemma-same-cohomology-adequate}
Soit $S$ un schéma. Soit $\mathcal{F}$ un
$\mathcal{O}$-module adéquat sur $(\Sch/S)_\tau$.
\begin{enumerate}
\item La restriction $\mathcal{F}|_{S_{Zar}}$ est un
$\mathcal{O}_S$-module quasi-cohérent sur le schéma $S$.
\item La restriction $\mathcal{F}|_{S_\etale}$ est le
module quasi-cohérent associé à $\mathcal{F}|_{S_{Zar}}$.
\item Pour tout schéma affine $U$ au-dessus de $S$, on a $H^q(U, \mathcal{F}) = 0$
pour tout $q > 0$.
\item Il existe un isomorphisme canonique
$$
H^q(S, \mathcal{F}|_{S_{Zar}}) =
H^q((\Sch/S)_\tau, \mathcal{F}).
$$
\end{enumerate}
\end{lemma}

\begin{proof}
D'après le
lemme \ref{lemma-adequate-flat}
et le
lemme \ref{lemma-adequate-local},
pour tout morphisme plat de schémas affines $U \to V$ au-dessus de $S$,
on a
$\mathcal{F}(U) = \mathcal{F}(V) \otimes_{\mathcal{O}(V)} \mathcal{O}(U)$.
Cela vaut en particulier si $U \subset V \subset S$ sont des ouverts affines de
$S$ ; ainsi $\mathcal{F}|_{S_{Zar}}$ est quasi-cohérent.
On obtient donc (1).

\medskip\noindent
Soit $S' \to S$ un morphisme étale de schémas.
Pour un ouvert affine $U \subset S'$ dont l'image est contenue dans un ouvert affine
$V \subset S$, on voit alors que
$\mathcal{F}(U) = \mathcal{F}(V) \otimes_{\mathcal{O}(V)} \mathcal{O}(U)$,
car $U \to V$ est étale, donc plat. Par conséquent,
$\mathcal{F}|_{S'_{Zar}}$ est l'image inverse de $\mathcal{F}|_{S_{Zar}}$.
Cela démontre (2).

\medskip\noindent
Nous allons appliquer
Cohomologie sur les sites,
lemme \ref{sites-cohomology-lemma-cech-vanish-collection},
au site $(\Sch/S)_\tau$, avec
$\mathcal{B}$ l'ensemble des schémas affines au-dessus de $S$ et
$\text{Cov}$ l'ensemble des $\tau$-recouvrements affines standard.
L'hypothèse (3) du lemme est satisfaite d'après
Descente, lemme \ref{descent-lemma-standard-covering-Cech},
et le
lemme \ref{lemma-adequate-fpqc}
pour le cas d'un recouvrement par un seul schéma affine.
On en conclut que $H^p(U, \mathcal{F}) = 0$ pour tout
schéma affine $U$ au-dessus de $S$. Cela démontre (3).
Exactement comme dans la démonstration de
Descente, proposition \ref{descent-proposition-same-cohomology-quasi-coherent},
cela entraîne l'égalité des cohomologies (4).
\end{proof}

\begin{remark}
\label{remark-compare}
Soit $S$ un schéma. On dispose de foncteurs
$u : \QCoh(\mathcal{O}_S) \to \textit{Adeq}(\mathcal{O})$
et
$v : \textit{Adeq}(\mathcal{O}) \to \QCoh(\mathcal{O}_S)$.
Le foncteur $u : \mathcal{F} \mapsto \mathcal{F}^a$
consiste à prendre le $\mathcal{O}$-module associé, qui est
adéquat d'après le
lemme \ref{lemma-adequate-characterize}.
Réciproquement, le foncteur $v$ est donné par restriction,
$v : \mathcal{G} \mapsto \mathcal{G}|_{S_{Zar}}$ ; voir le
lemme \ref{lemma-same-cohomology-adequate}.
Puisque $\mathcal{F}^a$ peut être décrit comme l'image inverse de
$\mathcal{F}$ par un morphisme de topos annelés
$((\Sch/S)_\tau, \mathcal{O}) \to (S_{Zar}, \mathcal{O}_S)$, voir
Descente, remarque \ref{descent-remark-change-topologies-ringed-sites},
et puisque la restriction est l'image directe, on voit que $u$ et $v$
sont adjoints comme suit
$$
\SheafHom_{\mathcal{O}_S}(\mathcal{F}, v\mathcal{G})
=
\SheafHom_\mathcal{O}(u\mathcal{F}, \mathcal{G})
$$
où $\mathcal{O}$ désigne le faisceau structural sur le grand site.
La description montre immédiatement que le morphisme d'adjonction
$\mathcal{F} \to vu\mathcal{F}$ est un isomorphisme pour tout faisceau
quasi-cohérent.
\end{remark}

\begin{lemma}
\label{lemma-sheafification-adequate}
{\emergencystretch=1em
Soit $S$ un schéma. Soit $\mathcal{F}$ un préfaisceau de $\mathcal{O}$-modules
sur $(\Sch/S)_\tau$. Si, pour tout schéma affine
$\Spec(A)$ au-dessus de $S$, le foncteur $F_{\mathcal{F}, A}$ est
adéquat, alors le faisceau associé à $\mathcal{F}$ est un
$\mathcal{O}$-module adéquat.\par}
\end{lemma}

\begin{proof}
Soit $U = \Spec(A)$ un schéma affine au-dessus de $S$.
Posons $F = F_{\mathcal{F}, A}$.
Le faisceau associé est $\mathcal{F}^\# = (\mathcal{F}^+)^+$ ; voir
Sites, section \ref{sites-section-sheafification}.
Par construction,
$$
(\mathcal{F})^+(U) =
\colim_\mathcal{U} \check{H}^0(\mathcal{U}, \mathcal{F})
$$
où la limite inductive porte sur les recouvrements dans le site $(\Sch/S)_\tau$.
Puisque $U$ est affine, il suffit de prendre la limite sur les
$\tau$-recouvrements affines standard
$\mathcal{U} = \{U_i \to U\}_{i \in I} =
\{\Spec(A_i) \to \Spec(A)\}_{i \in I}$ de $U$.
Puisque chaque $A \to A_i$ et chaque $A \to A_i \otimes_A A_j$ est plat, on voit que
$$
\check{H}^0(\mathcal{U}, \mathcal{F}) =
\Ker(\prod F(A) \otimes_A A_i \to \prod F(A) \otimes_A A_i \otimes_A A_j)
$$
d'après le
lemme \ref{lemma-adequate-flat}.
Puisque $A \to \prod A_i$ est fidèlement plat, ce module
est toujours canoniquement isomorphe à $F(A)$ d'après
Descente, lemme \ref{descent-lemma-ff-exact}.
Ainsi le préfaisceau $(\mathcal{F})^+$ a la même valeur que
$\mathcal{F}$ sur tous les schémas affines au-dessus de $S$. En répétant une fois
l'argument, on en déduit la même assertion pour $\mathcal{F}^\# = ((\mathcal{F})^+)^+$.
Ainsi $F_{\mathcal{F}, A} = F_{\mathcal{F}^\#, A}$, et l'on conclut
que $\mathcal{F}^\#$ est adéquat.
\end{proof}

\begin{lemma}
\label{lemma-abelian-adequate}
Soit $S$ un schéma.
\begin{enumerate}
\item La catégorie $\textit{Adeq}(\mathcal{O})$ est abélienne.
\item Le foncteur
$\textit{Adeq}(\mathcal{O}) \to
\textit{Mod}((\Sch/S)_\tau, \mathcal{O})$
est exact.
\item Si $0 \to \mathcal{F}_1 \to \mathcal{F}_2 \to \mathcal{F}_3 \to 0$
est une suite exacte courte de $\mathcal{O}$-modules et si
$\mathcal{F}_1$ et $\mathcal{F}_3$ sont adéquats, alors
$\mathcal{F}_2$ est adéquat.
\item La catégorie $\textit{Adeq}(\mathcal{O})$ admet des limites inductives, et
le foncteur
$$
\begin{aligned}
\textit{Adeq}(\mathcal{O}) & \longrightarrow \\
& \textit{Mod}((\Sch/S)_\tau, \mathcal{O})
\end{aligned}
$$
commute à ces limites.
\end{enumerate}
\end{lemma}

\begin{proof}
Soit $\varphi : \mathcal{F} \to \mathcal{G}$ une application de
$\mathcal{O}$-modules adéquats. Pour démontrer (1) et (2), il suffit de montrer que
$\mathcal{K} = \Ker(\varphi)$ et
$\mathcal{Q} = \Coker(\varphi)$, calculés dans
$\textit{Mod}((\Sch/S)_\tau, \mathcal{O})$, sont adéquats.
Soit $U = \Spec(A)$ un schéma affine au-dessus de $S$.
Posons $F = F_{\mathcal{F}, A}$ et $G = F_{\mathcal{G}, A}$.
D'après les
lemmes \ref{lemma-kernel-adequate} et
\ref{lemma-cokernel-adequate},
le noyau $K$ et le conoyau $Q$ de l'application induite
$F \to G$ sont des foncteurs adéquats.
Comme le noyau se calcule au niveau des préfaisceaux, on voit
que $K = F_{\mathcal{K}, A}$, et l'on conclut que $\mathcal{K}$ est adéquat.
Pour démontrer le résultat pour le conoyau, notons $\mathcal{Q}'$ le conoyau
préfaisceautique de $\varphi$. Alors $Q = F_{\mathcal{Q}', A}$ et
$\mathcal{Q} = (\mathcal{Q}')^\#$. Ainsi $\mathcal{Q}$
est adéquat d'après le
lemme \ref{lemma-sheafification-adequate}.

\medskip\noindent
Soit $0 \to \mathcal{F}_1 \to \mathcal{F}_2 \to \mathcal{F}_3 \to 0$
une suite exacte courte de $\mathcal{O}$-modules, avec
$\mathcal{F}_1$ et $\mathcal{F}_3$ adéquats.
Soit $U = \Spec(A)$ un schéma affine au-dessus de $S$.
Posons $F_i = F_{\mathcal{F}_i, A}$. La suite de foncteurs
$$
0 \to F_1 \to F_2 \to F_3 \to 0
$$
est exacte, car, pour $V = \Spec(B)$ affine au-dessus de $U$, on a
$H^1(V, \mathcal{F}_1) = 0$ d'après le
lemme \ref{lemma-same-cohomology-adequate}.
Puisque $F_1$ et $F_3$ sont des foncteurs adéquats d'après le
lemme \ref{lemma-adequate-local},
on voit que $F_2$ est adéquat d'après le
lemme \ref{lemma-extension-adequate}.
Ainsi $\mathcal{F}_2$ est adéquat.

\medskip\noindent
Soit $\mathcal{I} \to \textit{Adeq}(\mathcal{O})$, $i \mapsto \mathcal{F}_i$,
un diagramme. Notons $\mathcal{F} = \colim_i \mathcal{F}_i$
la limite inductive calculée dans
$\textit{Mod}((\Sch/S)_\tau, \mathcal{O})$.
Pour démontrer (4), il suffit de montrer que $\mathcal{F}$ est adéquat.
Soit $\mathcal{F}' = \colim_i \mathcal{F}_i$ la limite inductive calculée
dans les préfaisceaux de $\mathcal{O}$-modules. Alors
$\mathcal{F} = (\mathcal{F}')^\#$.
Soit $U = \Spec(A)$ un schéma affine au-dessus de $S$.
Posons $F_i = F_{\mathcal{F}_i, A}$. D'après le
lemme \ref{lemma-colimit-adequate},
le foncteur $\colim_i F_i = F_{\mathcal{F}', A}$ est adéquat.
Le lemme \ref{lemma-sheafification-adequate}
montre que $\mathcal{F}$ est adéquat.
\end{proof}

\noindent
Le lemme suivant affirme que l'image directe totale
$Rf_*\mathcal{F}$ d'un module adéquat par un morphisme quasi-compact et
quasi-séparé est un complexe dont les faisceaux de cohomologie
sont adéquats.

\begin{lemma}
\label{lemma-direct-image-adequate}
Soit $f : T \to S$ un morphisme quasi-compact et quasi-séparé
de schémas. Pour tout $\mathcal{O}_T$-module adéquat sur
$(\Sch/T)_\tau$, l'image directe
$f_*\mathcal{F}$ et les images directes supérieures $R^if_*\mathcal{F}$
sont des $\mathcal{O}_S$-modules adéquats sur $(\Sch/S)_\tau$.
\end{lemma}

\begin{proof}
Expliquons d'abord comment calculer les images directes supérieures.
Choisissons une résolution injective $\mathcal{F} \to \mathcal{I}^\bullet$.
Alors $R^if_*\mathcal{F}$ est le $i$-ième faisceau de cohomologie du
complexe $f_*\mathcal{I}^\bullet$.
Ainsi $R^if_*\mathcal{F}$ est le faisceau associé au préfaisceau
qui associe à un objet $U/S$ de $(\Sch/S)_\tau$
le module
\begin{align*}
\frac{\Ker(f_*\mathcal{I}^i(U) \to f_*\mathcal{I}^{i + 1}(U))}
{\Im(f_*\mathcal{I}^{i - 1}(U) \to f_*\mathcal{I}^i(U))}
& =
\frac{\Ker(\mathcal{I}^i(U \times_S T) \to
\mathcal{I}^{i + 1}(U \times_S T))}
{\Im(\mathcal{I}^{i - 1}(U \times_S T) \to \mathcal{I}^i(U \times_S T))}
\\
& =
H^i(U \times_S T, \mathcal{F}) \\
& = H^i((\Sch/U \times_S T)_\tau,
\mathcal{F}|_{(\Sch/U \times_S T)_\tau}) \\
& = H^i(U \times_S T, \mathcal{F}|_{(U \times_S T)_{Zar}})
\end{align*}
La première égalité résulte de
Topologies, lemme \ref{topologies-lemma-morphism-big-fppf}
(et de ses analogues pour les autres topologies),
la deuxième de la définition de la cohomologie de $\mathcal{F}$
au-dessus d'un objet de $(\Sch/T)_\tau$,
la troisième de
Cohomologie sur les sites, lemme \ref{sites-cohomology-lemma-cohomology-of-open},
et la dernière du
lemme \ref{lemma-same-cohomology-adequate}.
Ainsi, d'après le
lemme \ref{lemma-sheafification-adequate},
il suffit de démontrer l'assertion énoncée dans le paragraphe suivant.

\medskip\noindent
Soit $A$ un anneau. Soit $T$ un schéma quasi-compact et quasi-séparé
sur $A$. Soit $\mathcal{F}$ un $\mathcal{O}_T$-module adéquat sur
$(\Sch/T)_\tau$. Pour une $A$-algèbre $B$, posons
$T_B = T \times_{\Spec(A)} \Spec(B)$ et notons
$\mathcal{F}_B = \mathcal{F}|_{(T_B)_{Zar}}$ la restriction de
$\mathcal{F}$ au petit site de Zariski de $T_B$.
(Rappelons qu'il s'agit d'un faisceau quasi-cohérent « usuel » sur le schéma
$T_B$ ; voir le
lemme \ref{lemma-same-cohomology-adequate}.)
Assertion : le foncteur
$$
B \longmapsto H^q(T_B, \mathcal{F}_B)
$$
est adéquat. Nous démontrerons le lemme par la méthode usuelle
consistant à découper $T$ en morceaux.

\medskip\noindent
Cas I : $T$ est affine. Dans ce cas, les schémas $T_B$ sont tous affines
et $H^q(T_B, \mathcal{F}_B) = 0$ pour tout $q \geq 1$.
Le foncteur $B \mapsto H^0(T_B, \mathcal{F}_B)$ est adéquat d'après le
lemme \ref{lemma-pushforward-adequate}.

\medskip\noindent
Cas II : $T$ est séparé. Soit $n$ le nombre minimal d'ouverts affines nécessaires
pour recouvrir $T$. Raisonnons par récurrence sur $n$. Le cas initial est le cas I.
Choisissons un recouvrement ouvert affine $T = V_1 \cup \ldots \cup V_n$.
Posons $V = V_1 \cup \ldots \cup V_{n - 1}$ et $U = V_n$. Observons que
$$
U \cap V = (V_1 \cap V_n) \cup \ldots \cup (V_{n - 1} \cap V_n)
$$
est également une réunion de $n - 1$ ouverts affines, puisque $T$ est séparé ; voir
Schémas, lemme \ref{schemes-lemma-characterize-separated}.
Notons que, pour chaque $B$, les changements de base $U_B$, $V_B$ et
$(U \cap V)_B = U_B \cap V_B$ se comportent de la même manière. Ainsi,
pour chaque $B$, on a une suite exacte longue
$$
\begin{aligned}
0 &\to H^0(T_B, \mathcal{F}_B) \to
H^0(U_B, \mathcal{F}_B) \oplus H^0(V_B, \mathcal{F}_B) \\
&\to H^0((U \cap V)_B, \mathcal{F}_B) \to
H^1(T_B, \mathcal{F}_B) \to \ldots
\end{aligned}
$$
fonctorielle en $B$ ; voir
Cohomologie, lemme \ref{cohomology-lemma-mayer-vietoris}.
Par l'hypothèse de récurrence, les foncteurs
$B \mapsto H^q(U_B, \mathcal{F}_B)$,
$B \mapsto H^q(V_B, \mathcal{F}_B)$ et
$B \mapsto H^q((U \cap V)_B, \mathcal{F}_B)$
sont adéquats. Les
lemmes \ref{lemma-kernel-adequate} et
\ref{lemma-cokernel-adequate}
montrent que notre foncteur $B \mapsto H^q(T_B, \mathcal{F}_B)$ est le
terme médian d'une suite exacte courte dont les termes extrêmes sont adéquats.
L'assertion résulte donc du
lemme \ref{lemma-extension-adequate}.

\medskip\noindent
Cas III : cas général quasi-compact et quasi-séparé.
La démonstration se fait encore par récurrence sur le nombre $n$ d'ouverts affines
nécessaires pour recouvrir $T$. Le cas initial $n = 1$ est le cas I.
Choisissons un recouvrement ouvert affine $T = V_1 \cup \ldots \cup V_n$.
Posons $V = V_1 \cup \ldots \cup V_{n - 1}$ et $U = V_n$. Notons que,
puisque $T$ est quasi-séparé, $U \cap V$ est un ouvert quasi-compact d'un
schéma affine ; le cas II s'y applique donc. La suite de l'argument
se déroule exactement comme dans le paragraphe précédent, et nous
l'omettons.
\end{proof}







\section{Modules adéquats parasites}
\label{section-parasitic}

\noindent
Dans cette section, nous commençons à comparer les modules adéquats et les modules
quasi-cohérents sur un schéma $S$. Rappelons qu'il existe des foncteurs
$u : \QCoh(\mathcal{O}_S) \to \textit{Adeq}(\mathcal{O})$
et
$v : \textit{Adeq}(\mathcal{O}) \to \QCoh(\mathcal{O}_S)$
satisfaisant à l'adjonction
$$
\SheafHom_{\QCoh(\mathcal{O}_S)}(\mathcal{F}, v\mathcal{G})
=
\SheafHom_{\textit{Adeq}(\mathcal{O})}(u\mathcal{F}, \mathcal{G})
$$
et tels que $\mathcal{F} \to vu\mathcal{F}$ soit un isomorphisme pour
tout faisceau quasi-cohérent $\mathcal{F}$ ; voir la
remarque \ref{remark-compare}.
Ainsi $u$ est un plongement pleinement fidèle, et l'on peut identifier
$\QCoh(\mathcal{O}_S)$ à une sous-catégorie pleine de
$\textit{Adeq}(\mathcal{O})$.
Le foncteur $v$ est exact, mais $u$ n'est pas exact à gauche en général.
Le noyau de $v$ est la sous-catégorie des modules adéquats parasites.

\medskip\noindent
Dans Descente, définition \ref{descent-definition-parasitic},
nous donnons la définition d'un module parasite.
Pour les modules adéquats, cette notion ne dépend pas
de la topologie choisie.

\begin{lemma}
\label{lemma-parasitic-adequate}
Soit $S$ un schéma.
Soit $\mathcal{F}$ un $\mathcal{O}$-module adéquat sur
$(\Sch/S)_\tau$. Les conditions suivantes sont équivalentes :
\begin{enumerate}
\item $v\mathcal{F} = 0$,
\item $\mathcal{F}$ est parasite,
\item $\mathcal{F}$ est parasite pour la topologie $\tau$,
\item $\mathcal{F}(U) = 0$ pour tout ouvert $U \subset S$, et
\item il existe un recouvrement ouvert affine $S = \bigcup U_i$
tel que $\mathcal{F}(U_i) = 0$ pour tout $i$.
\end{enumerate}
\end{lemma}

\begin{proof}
Les implications (2) $\Rightarrow$ (3) $\Rightarrow$ (4) $\Rightarrow$ (5)
résultent immédiatement des définitions. Supposons (5). Supposons que
$S = \bigcup U_i$ soit un recouvrement ouvert affine tel que $\mathcal{F}(U_i) = 0$
pour tout $i$. Soit $V \to S$ un morphisme plat. Il existe un recouvrement ouvert
affine $V = \bigcup V_j$ tel que chaque $V_j$ ait son image dans un
$U_i$. Puisque le morphisme $V_j \to S$ est plat, $V_j \to U_i$ l'est aussi.
L'homomorphisme d'anneaux correspondant
$A_i = \mathcal{O}(U_i) \to \mathcal{O}(V_j) = B_j$ est donc plat. Ainsi, d'après le
lemme \ref{lemma-adequate-local}
et le
lemme \ref{lemma-adequate-flat},
on voit que $\mathcal{F}(U_i) \otimes_{A_i} B_j \to \mathcal{F}(V_j)$
est un isomorphisme. Ainsi $\mathcal{F}(V_j) = 0$. Puisque $\mathcal{F}$ est
un faisceau pour la topologie de Zariski, on conclut que $\mathcal{F}(V) = 0$.
On voit de cette manière que (5) entraîne (2).

\medskip\noindent
Cela démontre l'équivalence de (2), (3), (4) et (5).
Comme (1) est équivalente à (3) (voir la
remarque \ref{remark-compare}),
on conclut que les cinq conditions sont équivalentes.
\end{proof}

\noindent
Soit $S$ un schéma.
La sous-catégorie des modules adéquats parasites est une sous-catégorie de Serre de
$\textit{Adeq}(\mathcal{O})$. Le quotient est la catégorie des
modules quasi-cohérents.

\begin{lemma}
\label{lemma-adequate-by-parasitic}
Soit $S$ un schéma. La sous-catégorie
$\mathcal{C} \subset \textit{Adeq}(\mathcal{O})$ des modules adéquats
parasites est une sous-catégorie de Serre. De plus, le foncteur $v$ induit
une équivalence de catégories
$$
\textit{Adeq}(\mathcal{O}) / \mathcal{C} = \QCoh(\mathcal{O}_S).
$$
\end{lemma}

\begin{proof}
La catégorie $\mathcal{C}$ est le noyau du foncteur exact
$v : \textit{Adeq}(\mathcal{O}) \to \QCoh(\mathcal{O}_S)$ ; voir le
lemme \ref{lemma-parasitic-adequate}.
C'est donc une sous-catégorie de Serre d'après
Homologie, lemme \ref{homology-lemma-kernel-exact-functor}.
D'après
Homologie, lemme \ref{homology-lemma-serre-subcategory-is-kernel},
on obtient un foncteur exact induit
$\overline{v} :
\textit{Adeq}(\mathcal{O}) / \mathcal{C}
\to
\QCoh(\mathcal{O}_S)$.
Puisque $u$ est un inverse à droite de $v$, on voit immédiatement que
$\overline{v}$ est essentiellement surjectif.
On voit que $\overline{v}$ est fidèle d'après
Homologie, lemme \ref{homology-lemma-quotient-by-kernel-exact-functor}.
Puisque $u$ est un inverse à droite de $v$, on conclut enfin que
$\overline{v}$ est pleinement fidèle.
\end{proof}

\begin{lemma}
\label{lemma-direct-image-parasitic-adequate}
Soit $f : T \to S$ un morphisme quasi-compact et quasi-séparé
de schémas. Pour tout $\mathcal{O}_T$-module adéquat parasite sur
$(\Sch/T)_\tau$, l'image directe
$f_*\mathcal{F}$ et les images directes supérieures $R^if_*\mathcal{F}$
sont des $\mathcal{O}_S$-modules adéquats parasites sur $(\Sch/S)_\tau$.
\end{lemma}

\begin{proof}
Nous avons déjà vu dans le
lemme \ref{lemma-direct-image-adequate}
que ces images directes supérieures sont adéquates.
Il suffit donc de montrer que
$(R^if_*\mathcal{F})(U_i) = 0$ pour tout $\tau$-recouvrement
ouvert $\{U_i \to S\}$. Or $R^if_*\mathcal{F}$
est parasite d'après
Descente, lemme \ref{descent-lemma-direct-image-parasitic}.
\end{proof}













\section{Catégories dérivées de modules adéquats, I}
\label{section-comparison}


\noindent
Soit $S$ un schéma. Nous poursuivons la discussion commencée dans la
section \ref{section-parasitic}.
Le foncteur exact $v$ induit un foncteur
$$
D(\textit{Adeq}(\mathcal{O}))
\longrightarrow
D(\QCoh(\mathcal{O}_S))
$$
et de même pour les variantes bornées.

\begin{lemma}
\label{lemma-quotient-easy}
Soit $S$ un schéma. Notons
$\mathcal{C} \subset \textit{Adeq}(\mathcal{O})$ la
sous-catégorie pleine constituée des modules adéquats parasites.
Alors
$$
D(\textit{Adeq}(\mathcal{O}))/D_\mathcal{C}(\textit{Adeq}(\mathcal{O}))
= D(\QCoh(\mathcal{O}_S))
$$
et de même pour les variantes bornées.
\end{lemma}

\begin{proof}
{\emergencystretch=1em
Cela résulte immédiatement de
Catégories dérivées, lemme \ref{derived-lemma-quotient-by-serre-easy}.\par}
\end{proof}

\noindent
Cherchons ensuite une description en sens inverse en considérant
les foncteurs
$$
K^+(\QCoh(\mathcal{O}_S))
\longrightarrow
K^+(\textit{Adeq}(\mathcal{O}))
\longrightarrow
D^+(\textit{Adeq}(\mathcal{O}))
\longrightarrow
D^+(\QCoh(\mathcal{O}_S)).
$$
Dans certains cas, la catégorie dérivée des modules adéquats est une localisation
de la catégorie homotopique des complexes de modules quasi-cohérents par rapport
aux quasi-isomorphismes universels. Soit $S$ un schéma. Une application de complexes
$\varphi : \mathcal{F}^\bullet \to \mathcal{G}^\bullet$
de $\mathcal{O}_S$-modules quasi-cohérents est appelée
{\it quasi-isomorphisme universel} si, pour tout morphisme de schémas
$f : T \to S$, l'image inverse $f^*\varphi$ est un quasi-isomorphisme.

\begin{lemma}
\label{lemma-describe-Dplus-adequate}
Soit $U = \Spec(A)$ un schéma affine.
La catégorie dérivée bornée inférieurement
$D^+(\textit{Adeq}(\mathcal{O}))$ est la localisation
de $K^+(\QCoh(\mathcal{O}_U))$ par rapport au système multiplicatif des
quasi-isomorphismes universels.
\end{lemma}

\begin{proof}
{\emergencystretch=2em
Si $\varphi : \mathcal{F}^\bullet \to \mathcal{G}^\bullet$
est un morphisme de complexes de
$\mathcal{O}_U$-modules quasi-cohérents, alors $u\varphi : u\mathcal{F}^\bullet \to
u\mathcal{G}^\bullet$ est un quasi-isomorphisme si et seulement si $\varphi$ est
un quasi-isomorphisme universel. Ainsi la collection $S$
des quasi-isomorphismes universels est un système multiplicatif
saturé compatible avec la structure triangulée d'après
Catégories dérivées, lemme \ref{derived-lemma-triangle-functor-localize}.
Ainsi $S^{-1}K^+(\QCoh(\mathcal{O}_U))$ existe et est une
catégorie triangulée ; voir
Catégories dérivées, proposition
\ref{derived-proposition-construct-localization}.
On obtient un foncteur canonique
$$
\begin{aligned}
can : S^{-1}K^+(\QCoh(\mathcal{O}_U)) &\longrightarrow \\
&D^{+}(\textit{Adeq}(\mathcal{O}))
\end{aligned}
$$
d'après
Catégories dérivées, lemme \ref{derived-lemma-universal-property-localization}.
\par}

\medskip\noindent
Notons que, presque par définition, tout module adéquat sur $U$ se plonge
dans un faisceau quasi-cohérent ; voir le
lemme \ref{lemma-adequate-characterize}. Ainsi, d'après
Catégories dérivées, lemme \ref{derived-lemma-subcategory-right-resolution},
étant donné $\mathcal{F}^\bullet \in \Ob(K^+(\textit{Adeq}(\mathcal{O})))$,
il existe un quasi-isomorphisme
$\mathcal{F}^\bullet \to u\mathcal{G}^\bullet$,
où $\mathcal{G}^\bullet \in \Ob(K^+(\QCoh(\mathcal{O}_U)))$.
Cela démontre que $can$ est essentiellement surjectif.

\medskip\noindent
De même, supposons que $\mathcal{F}^\bullet$ et $\mathcal{G}^\bullet$
soient des complexes bornés inférieurement de $\mathcal{O}_U$-modules quasi-cohérents.
Un morphisme dans $D^+(\textit{Adeq}(\mathcal{O}))$ entre ces
complexes est constitué d'un couple $f : u\mathcal{F}^\bullet \to \mathcal{H}^\bullet$
et $s : u\mathcal{G}^\bullet \to \mathcal{H}^\bullet$, où $s$
est un quasi-isomorphisme. Choisissons un quasi-isomorphisme
$s' : \mathcal{H}^\bullet \to u\mathcal{E}^\bullet$. On voit alors que
$s' \circ f : \mathcal{F} \to \mathcal{E}^\bullet$ et le
quasi-isomorphisme universel
$s' \circ s : \mathcal{G}^\bullet \to \mathcal{E}^\bullet$ donnent
un morphisme dans $S^{-1}K^{+}(\QCoh(\mathcal{O}_U))$ ayant pour image
le morphisme donné. Cela démontre la plénitude.
La fidélité se démontre de manière analogue.
\end{proof}

\begin{lemma}
\label{lemma-right-adjoint-adequate}
Soit $U = \Spec(A)$ un schéma affine.
Le foncteur d'inclusion
$$
\textit{Adeq}(\mathcal{O}) \to
\textit{Mod}((\Sch/U)_\tau, \mathcal{O})
$$
admet un adjoint à droite $A$\footnote{C'est l'« adéquateur ».}.
De plus, le morphisme d'adjonction
$A(\mathcal{F}) \to \mathcal{F}$ est un isomorphisme pour tout
module adéquat $\mathcal{F}$.
\end{lemma}

\begin{proof}
D'après
Topologies, lemme \ref{topologies-lemma-affine-big-site-fppf}
(et de même pour les autres topologies),
on peut travailler avec les $\mathcal{O}$-modules sur $(\textit{Aff}/U)_\tau$.
Notons $\mathcal{P}$ la catégorie des foncteurs à valeurs
dans les modules sur $\textit{Alg}_A$ et $\mathcal{A}$ la catégorie des foncteurs adéquats
sur $\textit{Alg}_A$. Notons $i : \mathcal{A} \to \mathcal{P}$
le foncteur d'inclusion. Notons $Q : \mathcal{P} \to \mathcal{A}$
la construction du lemme \ref{lemma-adjoint}.
On a le diagramme commutatif
\begin{equation}
\label{equation-categories}
\vcenter{
\xymatrix{
\textit{Adeq}(\mathcal{O}) \ar[r]_-k \ar@{=}[d] &
\textit{Mod}((\textit{Aff}/U)_\tau, \mathcal{O}) \ar[r]_-j &
\textit{PMod}((\textit{Aff}/U)_\tau, \mathcal{O}) \ar@{=}[d] \\
\mathcal{A} \ar[rr]^-i & & \mathcal{P}
}
}
\end{equation}
L'égalité verticale de gauche est le
lemme \ref{lemma-adequate-affine},
et l'égalité verticale de droite a été expliquée dans la
section \ref{section-quasi-coherent}.
Définissons $A(\mathcal{F}) = Q(j(\mathcal{F}))$.
Puisque $j$ est pleinement fidèle, il en résulte immédiatement que $A$
est adjoint à droite du foncteur d'inclusion $k$. De plus, puisque
$k$ est lui aussi pleinement fidèle, la dernière assertion en résulte formellement.
\end{proof}

\noindent
Le foncteur $A$ est un adjoint à droite, donc est exact à gauche. Puisque le foncteur
d'inclusion est exact, voir le
lemme \ref{lemma-abelian-adequate},
on conclut que $A$ transforme les objets injectifs en objets injectifs et que
la catégorie $\textit{Adeq}(\mathcal{O})$ possède suffisamment d'objets injectifs ; voir
Homologie, lemme \ref{homology-lemma-adjoint-enough-injectives},
et
Injectifs, théorème \ref{injectives-theorem-sheaves-modules-injectives}.
Cela résulte également de l'équivalence dans
(\ref{equation-categories})
et du
lemme \ref{lemma-enough-injectives}.

\begin{lemma}
\label{lemma-RA-zero}
Soit $U = \Spec(A)$ un schéma affine.
Pour tout objet $\mathcal{F}$ de $\textit{Adeq}(\mathcal{O})$,
on a $R^pA(\mathcal{F}) = 0$ pour tout $p > 0$, où $A$ est
comme dans le
lemme \ref{lemma-right-adjoint-adequate}.
\end{lemma}

\begin{proof}
Avec les notations de la démonstration du
lemme \ref{lemma-right-adjoint-adequate},
choisissons une résolution injective $k(\mathcal{F}) \to \mathcal{I}^\bullet$
dans la catégorie des $\mathcal{O}$-modules sur $(\textit{Aff}/U)_\tau$.
D'après
Cohomologie sur les sites, lemme \ref{sites-cohomology-lemma-include-modules},
et le
lemme \ref{lemma-same-cohomology-adequate},
le complexe $j(\mathcal{I}^\bullet)$ est exact.
D'autre part, chaque $j(\mathcal{I}^n)$ est un objet injectif
de la catégorie des préfaisceaux de modules d'après
Cohomologie sur les sites, lemme
\ref{sites-cohomology-lemma-injective-module-injective-presheaf}.
Il en résulte que $R^pA(\mathcal{F}) = R^pQ(j(k(\mathcal{F})))$.
Le résultat découle donc du
lemme \ref{lemma-RQ-zero}.
\end{proof}

\noindent
Soit $S$ un schéma. D'après la discussion de la
section \ref{section-adequate},
le plongement
$\textit{Adeq}(\mathcal{O}) \subset
\textit{Mod}((\Sch/S)_\tau, \mathcal{O})$
présente $\textit{Adeq}(\mathcal{O})$ comme une sous-catégorie de Serre faible de
la catégorie de tous les $\mathcal{O}$-modules. Notons
$$
D_{\textit{Adeq}}(\mathcal{O}) \subset
D(\mathcal{O}) = D(\textit{Mod}((\Sch/S)_\tau, \mathcal{O}))
$$
la sous-catégorie triangulée des complexes dont les faisceaux de cohomologie
sont adéquats ; voir
Catégories dérivées, section \ref{derived-section-triangulated-sub}.
On obtient un foncteur canonique
$$
D(\textit{Adeq}(\mathcal{O}))
\longrightarrow
D_{\textit{Adeq}}(\mathcal{O})
$$
voir
Catégories dérivées, équation (\ref{derived-equation-compare}).

\begin{lemma}
\label{lemma-bounded-below}
Si $U = \Spec(A)$ est un schéma affine, alors la variante
bornée inférieurement
\begin{equation}
\label{equation-compare-bounded-adequate}
D^+(\textit{Adeq}(\mathcal{O}))
\longrightarrow
D^+_{\textit{Adeq}}(\mathcal{O})
\end{equation}
du foncteur ci-dessus est une équivalence.
\end{lemma}

\begin{proof}
Soit $A : \textit{Mod}(\mathcal{O}) \to \textit{Adeq}(\mathcal{O})$
l'adjoint à droite du foncteur d'inclusion construit dans le
lemme \ref{lemma-right-adjoint-adequate}.
Puisque $A$ est exact à gauche et que $\textit{Mod}(\mathcal{O})$
possède suffisamment d'objets injectifs, $A$ admet un foncteur dérivé à droite
$RA : D^+_{\textit{Adeq}}(\mathcal{O}) \to D^+(\textit{Adeq}(\mathcal{O}))$.
Nous affirmons que $RA$ est un quasi-inverse de
(\ref{equation-compare-bounded-adequate}).
Le fait essentiel est que, si $\mathcal{F}$ est un module adéquat,
le morphisme d'adjonction $\mathcal{F} \to RA(\mathcal{F})$ est un
quasi-isomorphisme d'après le lemme \ref{lemma-RA-zero}.

\medskip\noindent
Pour démontrer complètement le lemme, il suffit en effet de montrer que :
\begin{enumerate}
\item étant donné $\mathcal{F}^\bullet \in K^+(\textit{Adeq}(\mathcal{O}))$,
l'application canonique $\mathcal{F}^\bullet \to RA(\mathcal{F}^\bullet)$
est un quasi-isomorphisme, et
\item étant donné $\mathcal{G}^\bullet \in K^+(\textit{Mod}(\mathcal{O}))$,
l'application canonique $RA(\mathcal{G}^\bullet) \to \mathcal{G}^\bullet$
est un quasi-isomorphisme.
\end{enumerate}
Les assertions (1) et (2) résultent toutes deux du fait essentiel par un argument
de suite spectrale utilisant l'une des suites spectrales de
Catégories dérivées, lemme \ref{derived-lemma-two-ss-complex-functor}.
Nous omettons certains détails.
\end{proof}

\begin{lemma}
\label{lemma-ext-adequate}
Soit $U = \Spec(A)$ un schéma affine.
Soient $\mathcal{F}$ et $\mathcal{G}$ des $\mathcal{O}$-modules adéquats.
Pour tout $i \geq 0$, l'application naturelle
$$
\Ext^i_{\textit{Adeq}(\mathcal{O})}(\mathcal{F}, \mathcal{G})
\longrightarrow
\Ext^i_{\textit{Mod}(\mathcal{O})}(\mathcal{F}, \mathcal{G})
$$
est un isomorphisme.
\end{lemma}

\begin{proof}
Par définition, ces groupes Ext se calculent comme ensembles de morphismes dans
la catégorie dérivée. Cela résulte donc immédiatement du
lemme \ref{lemma-bounded-below}.
\end{proof}









\section{Extensions pures}
\label{section-pure}

\noindent
Nous voulons caractériser en termes algébriques les extensions de faisceaux
quasi-cohérents sur le grand site d'un schéma affine. Pour cela,
nous introduisons la notion suivante.

\begin{definition}
\label{definition-pure}
Soit $A$ un anneau.
\begin{enumerate}
\item Un $A$-module $P$ est dit {\it projectif pur}
si, pour toute suite universellement exacte
$0 \to K \to M \to N \to 0$ de $A$-modules, la suite
$0 \to \Hom_A(P, K) \to \Hom_A(P, M) \to \Hom_A(P, N) \to 0$
est exacte.
\item Un $A$-module $I$ est dit {\it injectif pur}
si, pour toute suite universellement exacte
$0 \to K \to M \to N \to 0$ de $A$-modules, la suite
$0 \to \Hom_A(N, I) \to \Hom_A(M, I) \to \Hom_A(K, I) \to 0$
est exacte.
\end{enumerate}
\end{definition}

\noindent
Caractérisons les modules projectifs purs.

\begin{lemma}
\label{lemma-pure-projective}
Soit $A$ un anneau.
\begin{enumerate}
\item Un module est projectif pur si et seulement s'il est
facteur direct d'une somme directe de $A$-modules de présentation finie.
\item Pour tout module $M$, il existe une suite universellement exacte
$0 \to N \to P \to M \to 0$, avec $P$ projectif pur.
\end{enumerate}
\end{lemma}

\begin{proof}
Notons d'abord qu'un $A$-module de présentation finie est projectif pur d'après
Algèbre, théorème \ref{algebra-theorem-universally-exact-criteria}.
Ainsi un facteur direct d'une somme directe de $A$-modules de présentation finie
est bien projectif pur. Soit $M$ un $A$-module quelconque.
Écrivons $M = \colim_{i \in I} P_i$ comme limite inductive filtrante de
$A$-modules de présentation finie. Considérons la suite
$$
0 \to N \to \bigoplus P_i \to M \to 0.
$$
Pour tout $A$-module de présentation finie $P$, l'application
$$
\Hom_A(P, \bigoplus P_i) \longrightarrow \Hom_A(P, M)
$$
est surjective, car toute application $P \to M$ se factorise par un certain $P_i$.
Ainsi, d'après
Algèbre, théorème \ref{algebra-theorem-universally-exact-criteria},
cette suite est universellement exacte. Cela démontre (2).
Si maintenant $M$ est projectif pur, alors la suite est scindée et
l'on voit que $M$ est facteur direct de $\bigoplus P_i$.
\end{proof}

\noindent
Caractérisons les modules injectifs purs.

\begin{lemma}
\label{lemma-pure-injective}
{\emergencystretch=2em
Soit $A$ un anneau. Pour tout $A$-module $M$, posons
$M^\vee = \Hom_\mathbf{Z}(M, \mathbf{Q}/\mathbf{Z})$.
\par}
\begin{enumerate}
\item Pour tout $A$-module $M$, le $A$-module $M^\vee$ est injectif pur.
\item Un $A$-module $I$ est injectif pur si et seulement si l'application
$I \to (I^\vee)^\vee$ est scindée.
\item Pour tout module $M$, il existe une suite universellement exacte
$0 \to M \to I \to N \to 0$, avec $I$ injectif pur.
\end{enumerate}
\end{lemma}

\begin{proof}
Nous utiliserons les propriétés du foncteur $M \mapsto M^\vee$ données dans
Compléments d'algèbre, section \ref{more-algebra-section-injectives-modules},
sans le mentionner à nouveau. L'assertion (1) résulte de ce que
$\Hom_A(N, M^\vee) =
\Hom_\mathbf{Z}(N \otimes_A M, \mathbf{Q}/\mathbf{Z})$
et de ce que $\mathbf{Q}/\mathbf{Z}$ est injectif dans la catégorie des
groupes abéliens. Ainsi, si $I \to (I^\vee)^\vee$ est scindée,
$I$ est injectif pur. Nous affirmons que, pour tout $A$-module $M$,
l'application d'évaluation $ev : M \to (M^\vee)^\vee$ est universellement injective.
Pour le voir, notons que $ev^\vee : ((M^\vee)^\vee)^\vee \to M^\vee$
admet un inverse à droite, à savoir $ev' : M^\vee \to ((M^\vee)^\vee)^\vee$.
Alors, pour tout $A$-module $N$, l'application du foncteur exact et fidèle
${}^\vee$ à l'application $N \otimes_A M \to N \otimes_A (M^\vee)^\vee$
donne
$$
\Hom_A(N, ((M^\vee)^\vee)^\vee) =
\Big(N \otimes_A (M^\vee)^\vee\Big)^\vee
\to
\Big(N \otimes_A M\Big)^\vee =
\Hom_A(N, M^\vee)
$$
qui est surjective en raison de l'existence de l'inverse à droite. D'où l'assertion.
Celle-ci entraîne (3) et la nécessité de la condition de (2).
\end{proof}

\noindent
Avant de poursuivre, faisons l'observation suivante, que nous utiliserons
fréquemment dans la suite de cette section.

\begin{lemma}
\label{lemma-split-universally-exact-sequence}
Soit $A$ un anneau.
\begin{enumerate}
\item Soit $L \to M \to N$ une suite universellement exacte
de $A$-modules. Posons $K = \Im(M \to N)$.
Alors $K \to N$ est universellement injective.
\item Tout complexe universellement exact
peut être décomposé en suites exactes courtes universellement exactes.
\end{enumerate}
\end{lemma}

\begin{proof}
Démonstration de (1). Pour tout $A$-module $T$, la suite
$L \otimes_A T \to M \otimes_A T \to K \otimes_A T \to 0$ est exacte
par exactitude à droite de $\otimes$. Par hypothèse, la suite
$L \otimes_A T \to M \otimes_A T \to N \otimes_A T$ est exacte.
Ces deux faits montrent que $K \otimes_A T \to N \otimes_A T$ est injective.

\medskip\noindent
L'assertion (2) signifie ceci : supposons que $M^\bullet$ soit un complexe
universellement exact de $A$-modules. Posons $K^i = \Ker(d^i) \subset M^i$.
Alors les suites exactes courtes $0 \to K^i \to M^i \to K^{i + 1} \to 0$
sont universellement exactes. Cela résulte immédiatement de (1).
\end{proof}

\begin{definition}
\label{definition-pure-resolution}
Soit $A$ un anneau. Soit $M$ un $A$-module.
\begin{enumerate}
\item Une {\it résolution par des projectifs purs} $P_\bullet \to M$
est une suite universellement exacte
$$
\ldots \to P_1 \to P_0 \to M \to 0
$$
où chaque $P_i$ est projectif pur.
\item Une {\it résolution par des injectifs purs} $M \to I^\bullet$ est une suite
universellement exacte
$$
0 \to M \to I^0 \to I^1 \to \ldots
$$
où chaque $I^i$ est injectif pur.
\end{enumerate}
\end{definition}

\noindent
Ces résolutions satisfont aux propriétés d'unicité usuelles parmi la classe
de toutes les résolutions à gauche ou à droite universellement exactes.

\begin{lemma}
\label{lemma-pure-projective-resolutions}
Soit $A$ un anneau.
\begin{enumerate}
\item Tout $A$-module admet une résolution par des projectifs purs.
\end{enumerate}
Soit $M \to N$ une application de $A$-modules.
Soit $P_\bullet \to M$ une résolution par des projectifs purs et
soit $N_\bullet \to N$ une résolution universellement exacte.
\begin{enumerate}
\item[(2)] Il existe une application de complexes $P_\bullet \to N_\bullet$
induisant l'application donnée
$$
M = \Coker(P_1 \to P_0) \to \Coker(N_1 \to N_0) = N
$$
\item[(3)] deux applications $\alpha, \beta : P_\bullet \to N_\bullet$
induisant la même application $M \to N$ sont homotopes.
\end{enumerate}
\end{lemma}

\begin{proof}
L'assertion (1) résulte immédiatement du
lemme \ref{lemma-pure-projective}.
Avant de démontrer (2) et (3), notons que, d'après le
lemme \ref{lemma-split-universally-exact-sequence},
on peut décomposer le complexe universellement exact $N_\bullet \to N \to 0$
en suites exactes courtes universellement exactes $0 \to K_0 \to N_0 \to N \to 0$
et $0 \to K_i \to N_i \to K_{i - 1} \to 0$.

\medskip\noindent
Démonstration de (2). Puisque $P_0$ est projectif pur,
on peut trouver une application $P_0 \to N_0$ relevant l'application $P_0 \to M \to N$.
On obtient une application induite $P_1 \to F_0 \to N_0$ dont l'image est contenue dans $K_0$.
Puisque $P_1$ est projectif pur, on peut la relever
en une application $P_1 \to N_1$. Celle-ci induit à son tour une application
$P_2 \to P_1 \to N_1$ dont l'image dans
$N_0$ est nulle, c'est-à-dire dont l'image est contenue dans $K_1$. On peut donc la relever en une application
$P_2 \to N_2$. On répète le procédé.

\medskip\noindent
Démonstration de (3). Pour montrer que $\alpha, \beta$ sont homotopes, il suffit
de montrer que la différence $\gamma = \alpha - \beta$ est homotope
à zéro. Notons que l'image de $\gamma_0 : P_0 \to N_0$
est contenue dans $K_0$. On peut donc relever
$\gamma_0$ en une application $h_0 : P_0 \to N_1$. Considérons l'application
$\gamma_1' = \gamma_1 - h_0 \circ d_{P, 1} : P_1 \to N_1$.
D'après le choix de $h_0$, on voit que l'image de $\gamma_1'$
est contenue dans $K_1$. Puisque $P_1$ est projectif pur, on peut relever
$\gamma_1'$ en une application $h_1 : P_1 \to N_2$. À ce stade, on a
$\gamma_1 = h_0 \circ d_{F, 1} + d_{N, 2} \circ h_1$. On répète le procédé.
\end{proof}

\begin{lemma}
\label{lemma-pure-injective-resolutions}
Soit $A$ un anneau.
\begin{enumerate}
\item Tout $A$-module admet une résolution par des injectifs purs.
\end{enumerate}
Soit $M \to N$ une application de $A$-modules.
Soit $M \to M^\bullet$ une résolution universellement exacte et
soit $N \to I^\bullet$ une résolution par des injectifs purs.
\begin{enumerate}
\item[(2)] Il existe une application de complexes $M^\bullet \to I^\bullet$
induisant l'application donnée
$$
M = \Ker(M^0 \to M^1) \to \Ker(I^0 \to I^1) = N
$$
\item[(3)] deux applications $\alpha, \beta : M^\bullet \to I^\bullet$
induisant la même application $M \to N$ sont homotopes.
\end{enumerate}
\end{lemma}

\begin{proof}
Ce lemme est dual du
lemme \ref{lemma-pure-projective-resolutions}.
La démonstration est identique, sauf qu'il faut renverser toutes les flèches.
\end{proof}

\noindent
Les résultats précédents permettent de définir les groupes d'extensions pures
comme suit. Soient $A$ un anneau et $M$, $N$ des $A$-modules.
Choisissons une résolution par des injectifs purs $N \to I^\bullet$. D'après le
lemme \ref{lemma-pure-injective-resolutions},
le complexe
$$
\Hom_A(M, I^\bullet)
$$
est bien défini à homotopie près. Son $i$-ième module de cohomologie
est donc un invariant bien défini de $M$ et $N$.

\begin{definition}
\label{definition-pure-ext}
Soient $A$ un anneau et $M$, $N$ des $A$-modules.
Le $i$-ième {\it module d'extensions pures} $\text{Pext}^i_A(M, N)$
est le $i$-ième module de cohomologie du complexe
$\Hom_A(M, I^\bullet)$, où $I^\bullet$ est une résolution par des injectifs purs
de $N$.
\end{definition}

\noindent
Attention : il n'est pas vrai qu'une suite exacte de $A$-modules donne
lieu à une suite exacte longue de groupes d'extensions pures. (Il faut
pour cela une suite universellement exacte.)
Rassemblons quelques faits qui résultent immédiatement de ce qui précède.

\begin{lemma}
\label{lemma-facts-pext}
Soit $A$ un anneau.
\begin{enumerate}
\item $\text{Pext}^i_A(M, N) = 0$ pour $i > 0$ dès que $N$ est injectif pur,
\item $\text{Pext}^i_A(M, N) = 0$ pour $i > 0$ dès que $M$ est projectif pur,
en particulier si $M$ est un $A$-module de présentation finie,
\item {\emergencystretch=3em
$\text{Pext}^i_A(M, N)$ est aussi le $i$-ième module de cohomologie
du complexe $\Hom_A(P_\bullet, N)$, où $P_\bullet$
est une résolution par des projectifs purs de $M$.\par}
\end{enumerate}
\end{lemma}

\begin{proof}
Pour (3), considérons le complexe double
$$
A^{\bullet, \bullet} = \Hom_A(P_\bullet, I^\bullet)
$$
{\emergencystretch=3em
Chacune de ses lignes est exacte sauf en degré $0$, où sa cohomologie
est $\Hom_A(M, I^q)$. Chacune de ses colonnes est exacte sauf en degré $0$,
où sa cohomologie est $\Hom_A(P_p, N)$. Ainsi les deux suites spectrales
associées à ce complexe dans
Homologie, section \ref{homology-section-double-complex},
dégénèrent, ce qui donne l'égalité.\par}
\end{proof}





\section{Groupes Ext supérieurs des faisceaux quasi-cohérents sur le grand site}
\label{section-big}

\noindent
Il se trouve que le foncteur à valeurs dans les modules $\underline{I}$ associé à
un module injectif pur $I$ donne un objet injectif dans la
catégorie des foncteurs adéquats sur $\textit{Alg}_A$.
Attention : il n'est pas vrai qu'un module projectif pur donne
un objet projectif dans la catégorie des foncteurs adéquats. Nous disposons néanmoins
de nombreux objets projectifs, à savoir les foncteurs linéairement adéquats.

\begin{lemma}
\label{lemma-pure-injective-injective-adequate}
Soit $A$ un anneau.
Soit $\mathcal{A}$ la catégorie des foncteurs adéquats sur $\textit{Alg}_A$.
Les objets injectifs de $\mathcal{A}$ sont exactement les foncteurs
$\underline{I}$, où $I$ est un $A$-module injectif pur.
\end{lemma}

\begin{proof}
Soit $I$ un objet injectif de $\mathcal{A}$.
Choisissons un plongement $I \to \underline{M}$ pour un certain $A$-module $M$.
Comme $I$ est injectif, on voit que $\underline{M} = I \oplus F$ pour un certain
foncteur à valeurs dans les modules $F$. Alors $M = I(A) \oplus F(A)$, et il en résulte
que $I = \underline{I(A)}$. On voit donc que tout objet injectif
est de la forme $\underline{I}$ pour un certain $A$-module $I$.
Il est clair que le module $I$ doit être injectif pur,
puisque toute suite universellement exacte $0 \to M \to N \to L \to 0$
donne lieu à une suite exacte
$0 \to \underline{M} \to \underline{N} \to \underline{L} \to 0$
de $\mathcal{A}$.

\medskip\noindent
Enfin, supposons que $I$ soit un
$A$-module injectif pur. Choisissons un plongement $\underline{I} \to J$
dans un objet injectif de $\mathcal{A}$ (voir le
lemme \ref{lemma-enough-injectives}).
Nous avons vu ci-dessus que $J = \underline{I'}$
pour un certain $A$-module $I'$ injectif pur. Puisque
$\underline{I} \to \underline{I'}$ est injective,
l'application $I \to I'$ est universellement injective. Par l'hypothèse sur $I$,
elle est scindée. Ainsi $\underline{I}$ est facteur direct de $J = \underline{I'}$,
donc est un objet injectif de la catégorie $\mathcal{A}$.
\end{proof}

\noindent
Soit $U = \Spec(A)$ un schéma affine. Soit $M$ un $A$-module.
Nous utiliserons la notation $M^a$ pour désigner le faisceau quasi-cohérent
de $\mathcal{O}$-modules sur $(\Sch/U)_\tau$ associé
au faisceau quasi-cohérent $\widetilde{M}$ sur $U$.
Nous disposons maintenant de toutes les notations nécessaires pour formuler le lemme suivant.

\begin{lemma}
\label{lemma-big-ext}
Soit $U = \Spec(A)$ un schéma affine. Soient $M$, $N$ des $A$-modules.
Pour tout $i$, on a un isomorphisme canonique
$$
\Ext^i_{\textit{Mod}(\mathcal{O})}(M^a, N^a) = \text{Pext}^i_A(M, N)
$$
fonctoriel en $M$ et $N$.
\end{lemma}

\begin{proof}
Construisons une flèche canonique de droite à gauche. Si
$N \to I^\bullet$ est une résolution par des injectifs purs, alors
$M^a \to (I^\bullet)^a$ est un complexe exact de
$\mathcal{O}$-modules (adéquats). Ainsi tout élément de $\text{Pext}^i_A(M, N)$
donne lieu à une application $N^a \to M^a[i]$ dans $D(\mathcal{O})$, c'est-à-dire à
un élément du groupe de gauche.

\medskip\noindent
Pour démontrer que cette application est un isomorphisme, notons que l'on peut remplacer
$\Ext^i_{\textit{Mod}(\mathcal{O})}(M^a, N^a)$ par
$\Ext^i_{\textit{Adeq}(\mathcal{O})}(M^a, N^a)$ ; voir le
lemme \ref{lemma-ext-adequate}.
Soit $\mathcal{A}$ la catégorie des foncteurs adéquats
sur $\textit{Alg}_A$. Nous avons vu que $\mathcal{A}$ est
équivalente à $\textit{Adeq}(\mathcal{O})$ ; voir le
lemme \ref{lemma-adequate-affine}, ainsi que la démonstration du
lemme \ref{lemma-right-adjoint-adequate}.
Il suffit donc maintenant de démontrer que
$$
\Ext^i_\mathcal{A}(\underline{M}, \underline{N}) =
\text{Pext}^i_A(M, N)
$$
Or cela résulte du
lemme \ref{lemma-pure-injective-injective-adequate},
puisqu'une résolution par des injectifs purs $N \to I^\bullet$ correspond exactement
à une résolution injective de $\underline{N}$ dans $\mathcal{A}$.
\end{proof}







\section{Catégories dérivées de modules adéquats, II}
\label{section-derived-categories}

\noindent
Soit $S$ un schéma. Notons $\mathcal{O}_S$ le faisceau structural de $S$
et $\mathcal{O}$ le faisceau structural du grand site $(\Sch/S)_\tau$.
Dans
Descente, remarque \ref{descent-remark-change-topologies-ringed},
nous avons construit un morphisme de sites annelés
\begin{equation}
\label{equation-compare-big-small}
f :
((\Sch/S)_\tau, \mathcal{O})
\longrightarrow
(S_{Zar}, \mathcal{O}_S).
\end{equation}
Dans les sections précédentes, nous avons vu que le foncteur
$f_* : \textit{Mod}(\mathcal{O}) \to \textit{Mod}(\mathcal{O}_S)$
transforme les faisceaux adéquats en faisceaux quasi-cohérents,
induit un foncteur exact
$v : \textit{Adeq}(\mathcal{O}) \to \QCoh(\mathcal{O}_S)$, et
qu'en fait $f_* = v$ induit une équivalence
$\textit{Adeq}(\mathcal{O})/\mathcal{C} \to \QCoh(\mathcal{O}_S)$,
où $\mathcal{C}$ est la sous-catégorie des modules adéquats parasites.
De plus, le foncteur $f^*$ transforme les modules quasi-cohérents
en modules adéquats et induit un foncteur
$u : \QCoh(\mathcal{O}_S) \to \textit{Adeq}(\mathcal{O})$
qui est adjoint à gauche de $v$.

\medskip\noindent
Il existe une relation très analogue entre
$D_{\textit{Adeq}}(\mathcal{O})$ et $D_\QCoh(S)$.
Expliquons d'abord pourquoi la catégorie $D_{\textit{Adeq}}(\mathcal{O})$
est indépendante de la topologie choisie.

\begin{remark}
\label{remark-D-adeq-independence-topology}
{\emergencystretch=3em
Soit $S$ un schéma.
Soient $\tau, \tau' \in \{Zar,\allowbreak \etale,\allowbreak smooth,\allowbreak syntomic,\allowbreak fppf\}$.
Notons $\mathcal{O}_\tau$, resp.\ $\mathcal{O}_{\tau'}$,
le faisceau structural $\mathcal{O}$ considéré comme faisceau sur
$(\Sch/S)_\tau$, resp.\ $(\Sch/S)_{\tau'}$.
Alors $D_{\textit{Adeq}}(\mathcal{O}_\tau)$ et
$D_{\textit{Adeq}}(\mathcal{O}_{\tau'})$ sont canoniquement isomorphes.
Cela résulte de Cohomologie sur les sites, lemme
\ref{sites-cohomology-lemma-compare-topologies-derived-adequate-modules}.
Supposons en effet que $\tau$ soit plus fine que la topologie $\tau'$, posons
$\mathcal{C} = (\Sch/S)_{fppf}$ et notons $\mathcal{B}$ la collection
des schémas affines au-dessus de $S$. Nous avons vérifié ci-dessus les hypothèses (1) et (2).
L'hypothèse (3) est claire, et l'hypothèse (4) résulte du
lemme \ref{lemma-same-cohomology-adequate}.\par}
\end{remark}

\begin{remark}
\label{remark-D-adeq-and-D-QCoh}
Soit $S$ un schéma. Le morphisme $f$, voir
(\ref{equation-compare-big-small}), induit
des foncteurs adjoints
$Rf_* : D_{\textit{Adeq}}(\mathcal{O}) \to D_\QCoh(S)$
et
$Lf^* : D_\QCoh(S) \to D_{\textit{Adeq}}(\mathcal{O})$.
De plus, $Rf_* Lf^* \cong \text{id}_{D_\QCoh(S)}$.

\medskip\noindent
Esquissons la démonstration. D'après la
remarque \ref{remark-D-adeq-independence-topology},
on peut supposer que la topologie $\tau$ est la topologie de Zariski.
Nous utiliserons l'existence des foncteurs dérivés totaux non bornés
$Lf^*$ et $Rf_*$ sur les $\mathcal{O}$-modules et leur
adjonction ; voir
Cohomologie sur les sites, lemme \ref{sites-cohomology-lemma-adjoint}.
Dans ce cas, $f_*$ est simplement la restriction à la sous-catégorie
$S_{Zar}$ de $(\Sch/S)_{Zar}$. Il est donc clair que
$Rf_* = f_*$ induit
$Rf_* : D_{\textit{Adeq}}(\mathcal{O}) \to D_\QCoh(S)$.
Supposons que $\mathcal{G}^\bullet$ soit un objet de
$D_\QCoh(S)$. On peut choisir un système
$\mathcal{K}_1^\bullet \to \mathcal{K}_2^\bullet \to \ldots$
de complexes bornés supérieurement de $\mathcal{O}_S$-modules plats dont les
applications de transition sont des injections scindées terme à terme, et un diagramme
$$
\xymatrix{
\mathcal{K}_1^\bullet \ar[d] \ar[r] &
\mathcal{K}_2^\bullet \ar[d] \ar[r] & \ldots \\
\tau_{\leq 1}\mathcal{G}^\bullet \ar[r] &
\tau_{\leq 2}\mathcal{G}^\bullet \ar[r] & \ldots
}
$$
ayant les propriétés (1), (2), (3) énumérées dans
Catégories dérivées, lemme \ref{derived-lemma-special-direct-system},
où $\mathcal{P}$ est la collection des $\mathcal{O}_S$-modules plats.
Alors $Lf^*\mathcal{G}^\bullet$ se calcule par
$\colim f^*\mathcal{K}_n^\bullet$ ; voir
Cohomologie sur les sites, lemmes \ref{sites-cohomology-lemma-pullback-K-flat} et
\ref{sites-cohomology-lemma-derived-base-change}
(notons que nos sites ont suffisamment de points d'après
Cohomologie étale, lemme \ref{etale-cohomology-lemma-points-fppf}).
Il faut voir que $H^i(Lf^*\mathcal{G}^\bullet) =
\colim H^i(f^*\mathcal{K}_n^\bullet)$ est adéquat pour chaque $i$. D'après le
lemme \ref{lemma-abelian-adequate},
on conclut qu'il suffit de montrer que
chaque $H^i(f^*\mathcal{K}_n^\bullet)$ est adéquat.

\medskip\noindent
Le caractère adéquat de $H^i(f^*\mathcal{K}_n^\bullet)$ est local sur $S$ ;
on peut donc supposer que $S = \Spec(A)$ est affine. Puisque $S$ est affine,
$D_\QCoh(S) = D(\QCoh(\mathcal{O}_S))$ ; voir
la discussion dans
Catégories dérivées de schémas, section
\ref{perfect-section-derived-quasi-coherent}.
Il existe donc un quasi-isomorphisme
$\mathcal{F}^\bullet \to \mathcal{K}_n^\bullet$,
où $\mathcal{F}^\bullet$ est un complexe borné supérieurement de modules
quasi-cohérents plats.
Alors $f^*\mathcal{F}^\bullet \to f^*\mathcal{K}_n^\bullet$ est un
quasi-isomorphisme, et les faisceaux de cohomologie de
$f^*\mathcal{F}^\bullet$ sont adéquats.

\medskip\noindent
La dernière assertion,
$Rf_* Lf^* \cong \text{id}_{D_\QCoh(S)}$,
résulte de la description explicite des foncteurs ci-dessus.
(En termes simples : si $\mathcal{F}$ est quasi-cohérent et si $p > 0$, alors
$L_pf^*\mathcal{F}$ est un module adéquat parasite.)
\end{remark}


\begin{remark}
\label{remark-conclusion}
La remarque \ref{remark-D-adeq-and-D-QCoh}
ci-dessus entraîne que l'on a une équivalence de catégories dérivées
$$
D_{\textit{Adeq}}(\mathcal{O})/D_\mathcal{C}(\mathcal{O})
\longrightarrow
D_\QCoh(S)
$$
où $\mathcal{C}$ est la catégorie des modules adéquats parasites.
Il est en effet clair que $D_\mathcal{C}(\mathcal{O})$ est le noyau
de $Rf_*$, d'où un foncteur comme indiqué. Pour tout objet $X$ de
$D_{\textit{Adeq}}(\mathcal{O})$, l'application $Lf^*Rf_*X \to X$ a pour image
un quasi-isomorphisme dans $D_\QCoh(S)$ ; ainsi
$Lf^*Rf_*X \to X$ est un isomorphisme dans
$D_{\textit{Adeq}}(\mathcal{O})/D_\mathcal{C}(\mathcal{O})$.
Enfin, pour des objets $X, Y$ de $D_{\textit{Adeq}}(\mathcal{O})$,
l'application
$$
Rf_* :
\Hom_{D_{\textit{Adeq}}(\mathcal{O})/D_\mathcal{C}(\mathcal{O})}(X, Y)
\to
\Hom_{D_\QCoh(S)}(Rf_*X, Rf_*Y)
$$
est bijective, puisque $Lf^*$ en donne un inverse (d'après les remarques ci-dessus).
\end{remark}









\begin{multicols}{2}[\section{Autres chapitres}]
\noindent
Préliminaires
\begin{enumerate}
\item \hyperref[introduction-section-phantom]{Introduction}
\item \hyperref[conventions-section-phantom]{Conventions}
\item \hyperref[sets-section-phantom]{Théorie des ensembles}
\item \hyperref[categories-section-phantom]{Catégories}
\item \hyperref[topology-section-phantom]{Topologie}
\item \hyperref[sheaves-section-phantom]{Faisceaux sur les espaces}
\item \hyperref[sites-section-phantom]{Sites et faisceaux}
\item \hyperref[stacks-section-phantom]{Champs}
\item \hyperref[fields-section-phantom]{Corps}
\item \hyperref[algebra-section-phantom]{Algèbre commutative}
\item \hyperref[brauer-section-phantom]{Groupes de Brauer}
\item \hyperref[homology-section-phantom]{Algèbre homologique}
\item \hyperref[derived-section-phantom]{Catégories dérivées}
\item \hyperref[simplicial-section-phantom]{Méthodes simpliciales}
\item \hyperref[more-algebra-section-phantom]{Compléments d'algèbre}
\item \hyperref[smoothing-section-phantom]{Lissage des morphismes d'anneaux}
\item \hyperref[modules-section-phantom]{Faisceaux de modules}
\item \hyperref[sites-modules-section-phantom]{Modules sur les sites}
\item \hyperref[injectives-section-phantom]{Objets injectifs}
\item \hyperref[cohomology-section-phantom]{Cohomologie des faisceaux}
\item \hyperref[sites-cohomology-section-phantom]{Cohomologie sur les sites}
\item \hyperref[dga-section-phantom]{Algèbre différentielle graduée}
\item \hyperref[dpa-section-phantom]{Algèbre à puissances divisées}
\item \hyperref[sdga-section-phantom]{Faisceaux différentiels gradués}
\item \hyperref[hypercovering-section-phantom]{Hyperrecouvrements}
\end{enumerate}
Schémas
\begin{enumerate}
\setcounter{enumi}{25}
\item \hyperref[schemes-section-phantom]{Schémas}
\item \hyperref[constructions-section-phantom]{Constructions de schémas}
\item \hyperref[properties-section-phantom]{Propriétés des schémas}
\item \hyperref[morphisms-section-phantom]{Morphismes de schémas}
\item \hyperref[coherent-section-phantom]{Cohomologie des schémas}
\item \hyperref[divisors-section-phantom]{Diviseurs}
\item \hyperref[limits-section-phantom]{Limites de schémas}
\item \hyperref[varieties-section-phantom]{Variétés}
\item \hyperref[topologies-section-phantom]{Topologies sur les schémas}
\item \hyperref[descent-section-phantom]{Descente}
\item \hyperref[perfect-section-phantom]{Catégories dérivées de schémas}
\item \hyperref[more-morphisms-section-phantom]{Compléments sur les morphismes}
\item \hyperref[flat-section-phantom]{Compléments sur la platitude}
\item \hyperref[groupoids-section-phantom]{Schémas en groupoïdes}
\item \hyperref[more-groupoids-section-phantom]{Compléments sur les schémas en groupoïdes}
\item \hyperref[etale-section-phantom]{Morphismes étales de schémas}
\end{enumerate}
Thèmes de théorie des schémas
\begin{enumerate}
\setcounter{enumi}{41}
\item \hyperref[chow-section-phantom]{Homologie de Chow}
\item \hyperref[intersection-section-phantom]{Théorie de l'intersection}
\item \hyperref[pic-section-phantom]{Schémas de Picard des courbes}
\item \hyperref[weil-section-phantom]{Théories cohomologiques de Weil}
\item \hyperref[adequate-section-phantom]{Modules adéquats}
\item \hyperref[dualizing-section-phantom]{Complexes dualisants}
\item \hyperref[duality-section-phantom]{Dualité pour les schémas}
\item \hyperref[discriminant-section-phantom]{Discriminants et différentes}
\item \hyperref[derham-section-phantom]{Cohomologie de de Rham}
\item \hyperref[local-cohomology-section-phantom]{Cohomologie locale}
\item \hyperref[algebraization-section-phantom]{Géométrie algébrique et formelle}
\item \hyperref[curves-section-phantom]{Courbes algébriques}
\item \hyperref[resolve-section-phantom]{Résolution des surfaces}
\item \hyperref[models-section-phantom]{Réduction semi-stable}
\item \hyperref[functors-section-phantom]{Foncteurs et morphismes}
\item \hyperref[equiv-section-phantom]{Catégories dérivées de variétés}
\item \hyperref[pione-section-phantom]{Groupes fondamentaux de schémas}
\item \hyperref[etale-cohomology-section-phantom]{Cohomologie étale}
\item \hyperref[crystalline-section-phantom]{Cohomologie cristalline}
\item \hyperref[proetale-section-phantom]{Cohomologie pro-étale}
\item \hyperref[relative-cycles-section-phantom]{Cycles relatifs}
\item \hyperref[more-etale-section-phantom]{Compléments de cohomologie étale}
\item \hyperref[trace-section-phantom]{La formule des traces}
\end{enumerate}
Espaces algébriques
\begin{enumerate}
\setcounter{enumi}{64}
\item \hyperref[spaces-section-phantom]{Espaces algébriques}
\item \hyperref[spaces-properties-section-phantom]{Propriétés des espaces algébriques}
\item \hyperref[spaces-morphisms-section-phantom]{Morphismes d'espaces algébriques}
\item \hyperref[decent-spaces-section-phantom]{Espaces algébriques décents}
\item \hyperref[spaces-cohomology-section-phantom]{Cohomologie des espaces algébriques}
\item \hyperref[spaces-limits-section-phantom]{Limites d'espaces algébriques}
\item \hyperref[spaces-divisors-section-phantom]{Diviseurs sur les espaces algébriques}
\item \hyperref[spaces-over-fields-section-phantom]{Espaces algébriques sur des corps}
\item \hyperref[spaces-topologies-section-phantom]{Topologies sur les espaces algébriques}
\item \hyperref[spaces-descent-section-phantom]{Descente et espaces algébriques}
\item \hyperref[spaces-perfect-section-phantom]{Catégories dérivées d'espaces}
\item \hyperref[spaces-more-morphisms-section-phantom]{Compléments sur les morphismes d'espaces}
\item \hyperref[spaces-flat-section-phantom]{Platitude sur les espaces algébriques}
\item \hyperref[spaces-groupoids-section-phantom]{Groupoïdes dans les espaces algébriques}
\item \hyperref[spaces-more-groupoids-section-phantom]{Compléments sur les groupoïdes dans les espaces}
\item \hyperref[bootstrap-section-phantom]{Amorçage}
\item \hyperref[spaces-pushouts-section-phantom]{Sommes amalgamées d'espaces algébriques}
\end{enumerate}
Thèmes de géométrie
\begin{enumerate}
\setcounter{enumi}{81}
\item \hyperref[spaces-chow-section-phantom]{Groupes de Chow des espaces}
\item \hyperref[groupoids-quotients-section-phantom]{Quotients de groupoïdes}
\item \hyperref[spaces-more-cohomology-section-phantom]{Compléments sur la cohomologie des espaces}
\item \hyperref[spaces-simplicial-section-phantom]{Espaces simpliciaux}
\item \hyperref[spaces-duality-section-phantom]{Dualité pour les espaces}
\item \hyperref[formal-spaces-section-phantom]{Espaces algébriques formels}
\item \hyperref[restricted-section-phantom]{Algébrisation des espaces formels}
\item \hyperref[spaces-resolve-section-phantom]{Résolution des surfaces revisitée}
\end{enumerate}
Théorie des déformations
\begin{enumerate}
\setcounter{enumi}{89}
\item \hyperref[formal-defos-section-phantom]{Théorie formelle des déformations}
\item \hyperref[defos-section-phantom]{Théorie des déformations}
\item \hyperref[cotangent-section-phantom]{Le complexe cotangent}
\item \hyperref[examples-defos-section-phantom]{Problèmes de déformation}
\end{enumerate}
Champs algébriques
\begin{enumerate}
\setcounter{enumi}{93}
\item \hyperref[algebraic-section-phantom]{Champs algébriques}
\item \hyperref[examples-stacks-section-phantom]{Exemples de champs}
\item \hyperref[stacks-sheaves-section-phantom]{Faisceaux sur les champs algébriques}
\item \hyperref[criteria-section-phantom]{Critères de représentabilité}
\item \hyperref[artin-section-phantom]{Axiomes d'Artin}
\item \hyperref[quot-section-phantom]{Espaces de Quot et de Hilbert}
\item \hyperref[stacks-properties-section-phantom]{Propriétés des champs algébriques}
\item \hyperref[stacks-morphisms-section-phantom]{Morphismes de champs algébriques}
\item \hyperref[stacks-limits-section-phantom]{Limites de champs algébriques}
\item \hyperref[stacks-cohomology-section-phantom]{Cohomologie des champs algébriques}
\item \hyperref[stacks-perfect-section-phantom]{Catégories dérivées de champs}
\item \hyperref[stacks-introduction-section-phantom]{Introduction aux champs algébriques}
\item \hyperref[stacks-more-morphisms-section-phantom]{Compléments sur les morphismes de champs}
\item \hyperref[stacks-geometry-section-phantom]{La géométrie des champs}
\end{enumerate}
Thèmes de théorie des espaces de modules
\begin{enumerate}
\setcounter{enumi}{107}
\item \hyperref[moduli-section-phantom]{Champs de modules}
\item \hyperref[moduli-curves-section-phantom]{Espaces de modules de courbes}
\end{enumerate}
Divers
\begin{enumerate}
\setcounter{enumi}{109}
\item \hyperref[examples-section-phantom]{Exemples}
\item \hyperref[exercises-section-phantom]{Exercices}
\item \hyperref[guide-section-phantom]{Guide bibliographique}
\item \hyperref[desirables-section-phantom]{Desiderata}
\item \hyperref[coding-section-phantom]{Style de programmation}
\item \hyperref[obsolete-section-phantom]{Obsolète}
\item \hyperref[fdl-section-phantom]{Licence de documentation libre GNU}
\item \hyperref[index-section-phantom]{Index généré automatiquement}
\end{enumerate}
\end{multicols}


\bibliography{my}
\bibliographystyle{amsalpha}

\end{document}
